% Options for packages loaded elsewhere
% Options for packages loaded elsewhere
\PassOptionsToPackage{unicode}{hyperref}
\PassOptionsToPackage{hyphens}{url}
\PassOptionsToPackage{dvipsnames,svgnames,x11names}{xcolor}
%
\documentclass[
  spanish,
  letterpaper,
  DIV=11,
  numbers=noendperiod]{scrreprt}
\usepackage{xcolor}
\usepackage{amsmath,amssymb}
\setcounter{secnumdepth}{0}
\usepackage{iftex}
\ifPDFTeX
  \usepackage[T1]{fontenc}
  \usepackage[utf8]{inputenc}
  \usepackage{textcomp} % provide euro and other symbols
\else % if luatex or xetex
  \usepackage{unicode-math} % this also loads fontspec
  \defaultfontfeatures{Scale=MatchLowercase}
  \defaultfontfeatures[\rmfamily]{Ligatures=TeX,Scale=1}
\fi
\usepackage{lmodern}
\ifPDFTeX\else
  % xetex/luatex font selection
  \setmonofont[]{JuliaMono}
\fi
% Use upquote if available, for straight quotes in verbatim environments
\IfFileExists{upquote.sty}{\usepackage{upquote}}{}
\IfFileExists{microtype.sty}{% use microtype if available
  \usepackage[]{microtype}
  \UseMicrotypeSet[protrusion]{basicmath} % disable protrusion for tt fonts
}{}
\makeatletter
\@ifundefined{KOMAClassName}{% if non-KOMA class
  \IfFileExists{parskip.sty}{%
    \usepackage{parskip}
  }{% else
    \setlength{\parindent}{0pt}
    \setlength{\parskip}{6pt plus 2pt minus 1pt}}
}{% if KOMA class
  \KOMAoptions{parskip=half}}
\makeatother
% Make \paragraph and \subparagraph free-standing
\makeatletter
\ifx\paragraph\undefined\else
  \let\oldparagraph\paragraph
  \renewcommand{\paragraph}{
    \@ifstar
      \xxxParagraphStar
      \xxxParagraphNoStar
  }
  \newcommand{\xxxParagraphStar}[1]{\oldparagraph*{#1}\mbox{}}
  \newcommand{\xxxParagraphNoStar}[1]{\oldparagraph{#1}\mbox{}}
\fi
\ifx\subparagraph\undefined\else
  \let\oldsubparagraph\subparagraph
  \renewcommand{\subparagraph}{
    \@ifstar
      \xxxSubParagraphStar
      \xxxSubParagraphNoStar
  }
  \newcommand{\xxxSubParagraphStar}[1]{\oldsubparagraph*{#1}\mbox{}}
  \newcommand{\xxxSubParagraphNoStar}[1]{\oldsubparagraph{#1}\mbox{}}
\fi
\makeatother

\usepackage{color}
\usepackage{fancyvrb}
\newcommand{\VerbBar}{|}
\newcommand{\VERB}{\Verb[commandchars=\\\{\}]}
\DefineVerbatimEnvironment{Highlighting}{Verbatim}{commandchars=\\\{\}}
% Add ',fontsize=\small' for more characters per line
\usepackage{framed}
\definecolor{shadecolor}{RGB}{241,243,245}
\newenvironment{Shaded}{\begin{snugshade}}{\end{snugshade}}
\newcommand{\AlertTok}[1]{\textcolor[rgb]{0.68,0.00,0.00}{#1}}
\newcommand{\AnnotationTok}[1]{\textcolor[rgb]{0.37,0.37,0.37}{#1}}
\newcommand{\AttributeTok}[1]{\textcolor[rgb]{0.40,0.45,0.13}{#1}}
\newcommand{\BaseNTok}[1]{\textcolor[rgb]{0.68,0.00,0.00}{#1}}
\newcommand{\BuiltInTok}[1]{\textcolor[rgb]{0.00,0.23,0.31}{#1}}
\newcommand{\CharTok}[1]{\textcolor[rgb]{0.13,0.47,0.30}{#1}}
\newcommand{\CommentTok}[1]{\textcolor[rgb]{0.37,0.37,0.37}{#1}}
\newcommand{\CommentVarTok}[1]{\textcolor[rgb]{0.37,0.37,0.37}{\textit{#1}}}
\newcommand{\ConstantTok}[1]{\textcolor[rgb]{0.56,0.35,0.01}{#1}}
\newcommand{\ControlFlowTok}[1]{\textcolor[rgb]{0.00,0.23,0.31}{\textbf{#1}}}
\newcommand{\DataTypeTok}[1]{\textcolor[rgb]{0.68,0.00,0.00}{#1}}
\newcommand{\DecValTok}[1]{\textcolor[rgb]{0.68,0.00,0.00}{#1}}
\newcommand{\DocumentationTok}[1]{\textcolor[rgb]{0.37,0.37,0.37}{\textit{#1}}}
\newcommand{\ErrorTok}[1]{\textcolor[rgb]{0.68,0.00,0.00}{#1}}
\newcommand{\ExtensionTok}[1]{\textcolor[rgb]{0.00,0.23,0.31}{#1}}
\newcommand{\FloatTok}[1]{\textcolor[rgb]{0.68,0.00,0.00}{#1}}
\newcommand{\FunctionTok}[1]{\textcolor[rgb]{0.28,0.35,0.67}{#1}}
\newcommand{\ImportTok}[1]{\textcolor[rgb]{0.00,0.46,0.62}{#1}}
\newcommand{\InformationTok}[1]{\textcolor[rgb]{0.37,0.37,0.37}{#1}}
\newcommand{\KeywordTok}[1]{\textcolor[rgb]{0.00,0.23,0.31}{\textbf{#1}}}
\newcommand{\NormalTok}[1]{\textcolor[rgb]{0.00,0.23,0.31}{#1}}
\newcommand{\OperatorTok}[1]{\textcolor[rgb]{0.37,0.37,0.37}{#1}}
\newcommand{\OtherTok}[1]{\textcolor[rgb]{0.00,0.23,0.31}{#1}}
\newcommand{\PreprocessorTok}[1]{\textcolor[rgb]{0.68,0.00,0.00}{#1}}
\newcommand{\RegionMarkerTok}[1]{\textcolor[rgb]{0.00,0.23,0.31}{#1}}
\newcommand{\SpecialCharTok}[1]{\textcolor[rgb]{0.37,0.37,0.37}{#1}}
\newcommand{\SpecialStringTok}[1]{\textcolor[rgb]{0.13,0.47,0.30}{#1}}
\newcommand{\StringTok}[1]{\textcolor[rgb]{0.13,0.47,0.30}{#1}}
\newcommand{\VariableTok}[1]{\textcolor[rgb]{0.07,0.07,0.07}{#1}}
\newcommand{\VerbatimStringTok}[1]{\textcolor[rgb]{0.13,0.47,0.30}{#1}}
\newcommand{\WarningTok}[1]{\textcolor[rgb]{0.37,0.37,0.37}{\textit{#1}}}

\usepackage{longtable,booktabs,array}
\newcounter{none} % for unnumbered tables
\usepackage{calc} % for calculating minipage widths
% Correct order of tables after \paragraph or \subparagraph
\usepackage{etoolbox}
\makeatletter
\patchcmd\longtable{\par}{\if@noskipsec\mbox{}\fi\par}{}{}
\makeatother
% Allow footnotes in longtable head/foot
\IfFileExists{footnotehyper.sty}{\usepackage{footnotehyper}}{\usepackage{footnote}}
\makesavenoteenv{longtable}
\usepackage{graphicx}
\makeatletter
\newsavebox\pandoc@box
\newcommand*\pandocbounded[1]{% scales image to fit in text height/width
  \sbox\pandoc@box{#1}%
  \Gscale@div\@tempa{\textheight}{\dimexpr\ht\pandoc@box+\dp\pandoc@box\relax}%
  \Gscale@div\@tempb{\linewidth}{\wd\pandoc@box}%
  \ifdim\@tempb\p@<\@tempa\p@\let\@tempa\@tempb\fi% select the smaller of both
  \ifdim\@tempa\p@<\p@\scalebox{\@tempa}{\usebox\pandoc@box}%
  \else\usebox{\pandoc@box}%
  \fi%
}
% Set default figure placement to htbp
\def\fps@figure{htbp}
\makeatother



\ifLuaTeX
\usepackage[bidi=basic,shorthands=off]{babel}
\else
\usepackage[bidi=default,shorthands=off]{babel}
\fi
\ifLuaTeX
  \usepackage{selnolig} % disable illegal ligatures
\fi


\setlength{\emergencystretch}{3em} % prevent overfull lines

\providecommand{\tightlist}{%
  \setlength{\itemsep}{0pt}\setlength{\parskip}{0pt}}



 
\usepackage[style=alphabetic,giveninits=true]{biblatex}
\addbibresource{references.bib}


\usepackage{amsthm}
\usepackage{mathpartir}
\usepackage{tikz}
\usepackage[normalem]{ulem}
\usepackage{scrlayer-scrpage}

\automark{section}
\setkomafont{pagefoot}{\upshape}
\cfoot*{}
\ofoot*{\pagemark}

\makeatletter
\def\redsquiggly{\bgroup \markoverwith{\textcolor{red}{\lower3.5\p@\hbox{\sixly \char58}}}\ULon}
\def\brownsquiggly{\bgroup \markoverwith{\textcolor[HTML]{B8860B}{\lower3.5\p@\hbox{\sixly \char58}}}\ULon}
\def\bluesquiggly{\bgroup \markoverwith{\textcolor[HTML]{1E90FF}{\lower3.5\p@\hbox{\sixly \char58}}}\ULon}
\makeatother

\renewcommand{\NormalTok}[1]{\textcolor[HTML]{000000}{#1}}
\renewcommand{\KeywordTok}[1]{\textcolor[HTML]{0000FF}{#1}}
\renewcommand{\SpecialCharTok}[1]{}
\renewcommand{\ErrorTok}[1]{\redsquiggly{#1}}
\renewcommand{\WarningTok}[1]{\redsquiggly{\textcolor[HTML]{0000FF}{#1}}}
\renewcommand{\StringTok}[1]{\textcolor[HTML]{A52A2A}{#1}}
\renewcommand{\CommentTok}[1]{\textcolor[HTML]{008000}{#1}}
\renewcommand{\ConstantTok}[1]{\textcolor[HTML]{DC143C}{#1}}
\renewcommand{\AnnotationTok}[1]{\brownsquiggly{#1}}
\renewcommand{\AlertTok}[1]{\brownsquiggly{\textcolor[HTML]{0000FF}{#1}}}
\renewcommand{\OtherTok}[1]{\bluesquiggly{#1}}
\renewcommand{\DocumentationTok}[1]{\bluesquiggly{\textcolor[HTML]{0000FF}{#1}}}

\AtBeginDocument{%
  \let\oldgreater\textgreater
  \renewcommand{\textgreater}{\null\oldgreater}%
}

\pagenumbering{roman}
\KOMAoption{captions}{tableheading}
\makeatletter
\@ifpackageloaded{bookmark}{}{\usepackage{bookmark}}
\makeatother
\makeatletter
\@ifpackageloaded{caption}{}{\usepackage{caption}}
\AtBeginDocument{%
\ifdefined\contentsname
  \renewcommand*\contentsname{Tabla de contenidos}
\else
  \newcommand\contentsname{Tabla de contenidos}
\fi
\ifdefined\listfigurename
  \renewcommand*\listfigurename{Listado de Figuras}
\else
  \newcommand\listfigurename{Listado de Figuras}
\fi
\ifdefined\listtablename
  \renewcommand*\listtablename{Listado de Tablas}
\else
  \newcommand\listtablename{Listado de Tablas}
\fi
\ifdefined\figurename
  \renewcommand*\figurename{Figura}
\else
  \newcommand\figurename{Figura}
\fi
\ifdefined\tablename
  \renewcommand*\tablename{Tabla}
\else
  \newcommand\tablename{Tabla}
\fi
}
\@ifpackageloaded{float}{}{\usepackage{float}}
\floatstyle{ruled}
\@ifundefined{c@chapter}{\newfloat{codelisting}{h}{lop}}{\newfloat{codelisting}{h}{lop}[chapter]}
\floatname{codelisting}{Listado}
\newcommand*\listoflistings{\listof{codelisting}{Listado de Listados}}
\usepackage{amsthm}
\theoremstyle{remark}
\AtBeginDocument{\renewcommand*{\proofname}{Demostración}}
\newtheorem*{remark}{Observación}
\newtheorem*{solution}{Solución}
\newtheorem{refremark}{Observación}[chapter]
\newtheorem{refsolution}{Solución}[chapter]
\makeatother
\makeatletter
\makeatother
\makeatletter
\@ifpackageloaded{caption}{}{\usepackage{caption}}
\@ifpackageloaded{subcaption}{}{\usepackage{subcaption}}
\makeatother
\usepackage{bookmark}
\IfFileExists{xurl.sty}{\usepackage{xurl}}{} % add URL line breaks if available
\urlstyle{same}
% fallback for those not using the hyperref driver hyperxmp:
\makeatletter
\@ifundefined{xmpquote}{\newcommand{\xmpquote}[1]{#1}}{}
\makeatother
\hypersetup{
  pdftitle={Una prueba transicional de la confluencia de la reducción \textbackslash eta},
  pdfauthor={Maximiliano Onofre Martínez},
  pdflang={es},
  colorlinks=true,
  linkcolor={blue},
  filecolor={Maroon},
  citecolor={Blue},
  urlcolor={Blue},
  pdfcreator={LaTeX via pandoc}}


\title{Una prueba transicional de la confluencia de la reducción
\(\eta\)}
\author{Maximiliano Onofre Martínez}
\date{}
\begin{document}
\maketitle

\renewcommand*\contentsname{Tabla de contenidos}
{
\setcounter{tocdepth}{1}
\tableofcontents
}

\bookmarksetup{startatroot}

\chapter*{Introducción}\label{introducciuxf3n}
\addcontentsline{toc}{chapter}{Introducción}

\markboth{Introducción}{Introducción}

\pagenumbering{arabic}

Un asistente de pruebas es un programa que permite expresar definiciones
y enunciados matemáticos en un lenguaje formal, y construir
demostraciones que el sistema comprueba según sus reglas lógicas.

Miranda-Perea y González Huesca proponen las \emph{pruebas
transicionales} para conciliar las demostraciones habituales en papel
con las realizadas en un asistente de pruebas. Una prueba transicional
consta de una demostración matemática escrita y una secuencia de
instrucciones que demuestra el mismo enunciado en el asistente. Cada
paso del argumento escrito debe corresponder a una secuencia de
instrucciones que exprese el mismo razonamiento y mantenga el orden de
los pasos.

La definición requiere también un uso equilibrado de las herramientas
del asistente. Solo las tareas operativas más elementales pueden
aparecer en la formalización sin un paso correspondiente en el texto
\autocite{miranda2022}.

Para construir esta correspondencia, trabajamos con el texto y el
código. Al formalizar, podemos necesitar precisar algún paso del
argumento escrito. La redacción también puede ayudarnos a organizar el
código para que sea más fácil seguir ese argumento.

En este trabajo presentamos una prueba transicional de la confluencia de
la reducción \(\eta\) en el cálculo \(\lambda\), utilizando Lean 4 para
su componente formal.

La reducción \(\eta\) permite simplificar una expresión que describe una
función. Por ejemplo, \(\lambda x. f \, x\) recibe un argumento y se
limita a aplicar \(f\) a ese argumento, por lo que podemos escribir
simplemente \(f\). En general, la reducción tiene la forma \[
\lambda x. M \, x\longrightarrow_\eta M, \quad x\notin\operatorname{fv}(M).
\] Este reemplazo también puede realizarse dentro de aplicaciones y
abstracciones. Cuando hay varias partes que podemos reducir, surge la
pregunta de si el orden elegido afecta al resultado. Una situación
familiar aparece al calcular \((2+3)\cdot(4+1)\): podemos comenzar por
cualquiera de las dos sumas y obtener expresiones intermedias distintas,
pero en ambos casos llegamos a \(25\). En el cálculo \(\lambda\), la
confluencia garantiza que, si dos sucesiones de reducciones parten del
mismo término, podemos continuar desde los términos obtenidos hasta
alcanzar uno común.

La confluencia de \(\eta\) es un resultado conocido. Tomamos como
referencia la demostración presentada por Barendregt
\autocite{barendregt1984}. Entre sus antecedentes formales se encuentran
los desarrollos de Nipkow en Isabelle \autocite{nipkow2001} y de
Momigliano y Sassella en Beluga \autocite{momigliano2024}. Estos
trabajos estudian también las reducciones \(\beta\) y \(\beta\eta\);
aquí nos concentramos en \(\eta\). Nuestro propósito es desarrollar una
exposición en la que podamos seguir el argumento matemático y reconocer
cómo se expresa cada paso en Lean.

Para representar los términos utilizamos la representación locally
nameless, siguiendo el trabajo de Charguéraud
\autocite{chargueraud2012}. En ella, las variables libres conservan sus
nombres y las ligadas se representan mediante números que indican a qué
abstracción hacen referencia. Esta elección permite seguir utilizando
nombres en los enunciados y las demostraciones. A lo largo del trabajo
explicamos las operaciones y los lemas que necesitamos para razonar con
esta representación.

La presentación transicional se concentra en la demostración de la
confluencia. Para los resultados auxiliares presentamos las
demostraciones o explicaciones matemáticas y sus enunciados en Lean. Sus
demostraciones formales se encuentran en el repositorio que acompaña al
trabajo.\footnote{Disponible en
  \url{https://github.com/m-ow/bachelors-thesis}.}

A lo largo del trabajo utilizamos nociones básicas de relaciones y
lógica, así como demostraciones por inducción. Partimos también de
cierta familiaridad con la notación habitual del cálculo \(\lambda\). En
cuanto a Lean, no suponemos experiencia previa: explicamos cómo leer los
enunciados y las demostraciones escritos en el asistente, para que el
lector pueda seguir el código junto con el argumento matemático.

\section*{Organización del trabajo}\label{organizaciuxf3n-del-trabajo}
\addcontentsline{toc}{section}{Organización del trabajo}

\markright{Organización del trabajo}

El capítulo 1 presenta las clausuras de relaciones, la propiedad del
diamante y la confluencia. El capítulo 2 introduce los términos en la
representación \emph{locally nameless}, sus operaciones y la clausura
local. El capítulo 3 desarrolla las propiedades de la reducción \(\eta\)
necesarias para la demostración. El capítulo 4 presenta la prueba
transicional de la confluencia de esta reducción. Finalmente, en la
conclusión se explica la aportación del trabajo y se propone como
continuación elaborar pruebas transicionales de la confluencia de las
reducciones \(\beta\) y \(\beta\eta\).

\bookmarksetup{startatroot}

\chapter{Relaciones}\label{relaciones}

En teoría de conjuntos se entiende una relación binaria \(R\) sobre un
conjunto \(A\) como un conjunto de pares ordenados, de modo que
\(R(a,b)\) significa \((a,b) \in R\). Una alternativa es pensar en \(R\)
como una función que, al aplicarse a \(a\) y \(b\), produce una
proposición que expresa que ambos elementos están relacionados.
Adoptaremos esta presentación en Lean mediante la siguiente definición:

\begin{Shaded}
\begin{Highlighting}[]
\KeywordTok{def}\NormalTok{ relation (X : }\KeywordTok{Type}\NormalTok{) := X → X → }\KeywordTok{Prop}
\end{Highlighting}
\end{Shaded}

Así, si \texttt{(R\ :\ relation\ X)} y \texttt{(a\ b\ :\ X)}, la
expresión \texttt{(R\ a\ b)} representa la proposición que afirma que
\texttt{a} está relacionado con \texttt{b} mediante \texttt{R}.

Sea \texttt{R\ :\ relation\ X}. La relación \texttt{R} es
\emph{reflexiva} si todo término de tipo \texttt{X} está relacionado
consigo mismo, y es \emph{transitiva} si, para cualesquiera términos
\texttt{a}, \texttt{b} y \texttt{c} de tipo \texttt{X}, de
\texttt{R\ a\ b} y \texttt{R\ b\ c} se sigue \texttt{R\ a\ c}. Estas
propiedades pueden expresarse en Lean de la siguiente manera:

\begin{Shaded}
\begin{Highlighting}[]
\NormalTok{variable \{X : }\KeywordTok{Type}\NormalTok{\}}

\KeywordTok{def}\NormalTok{ reflexive (R : relation X) : }\KeywordTok{Prop}\NormalTok{ :=}
\NormalTok{  ∀ a : X, R a a}

\KeywordTok{def}\NormalTok{ transitive (R : relation X) : }\KeywordTok{Prop}\NormalTok{ :=}
\NormalTok{  ∀ a b c : X, R a b → R b c → R a c}
\end{Highlighting}
\end{Shaded}

Las llaves en \texttt{variable\ \{X\ :\ Type\}} indican que \texttt{X}
es un argumento \emph{implícito}: en lugar de exigir que se proporcione,
Lean lo infiere a partir del contexto. Así, como \texttt{R} tiene tipo
\texttt{relation\ X}, Lean determina \texttt{X} a partir del tipo de
\texttt{R}, y basta escribir \texttt{reflexive\ R} o
\texttt{transitive\ R} sin indicar el tipo sobre el cual está definida
la relación.

\section{1.1. Clausuras}\label{clausuras}

A continuación presentamos las tres clausuras de \(R\) que utilizaremos,
junto con sus reglas de inferencia y sus definiciones en Lean.
Interpretaremos \(R(a,b)\) como un paso de \(a\) a \(b\).

I. La clausura reflexiva \(R^{=}\) relaciona dos elementos si son
iguales o si están relacionados por \(R\). Se define mediante las
siguientes reglas:

\[
\inferrule*[right=\textsc{refl}]
  { }
  {R^=(a,a)}
\qquad
\inferrule*[right=\textsc{step}]
  {R(a,b)}
  {R^=(a,b)}
\]

\begin{Shaded}
\begin{Highlighting}[]
\NormalTok{inductive clos\_refl (R : relation X) : relation X}
\NormalTok{  | refl : ∀ a : X,}
\NormalTok{      clos\_refl R a a}
\NormalTok{  | step : ∀ a b : X,}
\NormalTok{      R a b → clos\_refl R a b}
\end{Highlighting}
\end{Shaded}

II. La clausura transitiva \(R^{+}\) relaciona dos elementos cuando
pueden conectarse mediante una sucesión finita de uno o más pasos de
\(R\). Sus reglas de inferencia son:

\[
\inferrule*[right=\textsc{step}]
  {R(a,b)}
  {R^+(a,b)}
\qquad
\inferrule*[right=\textsc{trans}]
  {R^+(a,b) \\
   R^+(b,c)}
  {R^+(a,c)}
\]

\begin{Shaded}
\begin{Highlighting}[]
\NormalTok{inductive clos\_trans (R : relation X) : relation X}
\NormalTok{  | step : ∀ a b : X,}
\NormalTok{      R a b → clos\_trans R a b}
\NormalTok{  | trans : ∀ a b c : X,}
\NormalTok{      clos\_trans R a b → clos\_trans R b c →}
\NormalTok{      clos\_trans R a c}
\end{Highlighting}
\end{Shaded}

III. La clausura reflexiva-transitiva \(R^*\) relaciona dos elementos
cuando pueden conectarse mediante una sucesión finita de cero o más
pasos de \(R\). Una sucesión de cero pasos comienza y termina en el
mismo elemento.

\[
\inferrule*[right=\textsc{refl}]
  { }
  {R^*(a,a)}
\qquad
\inferrule*[right=\textsc{step}]
  {R(a,b)}
  {R^*(a,b)}
\qquad
\inferrule*[right=\textsc{trans}]
  {R^*(a,b) \\
   R^*(b,c)}
  {R^*(a,c)}
\]

\begin{Shaded}
\begin{Highlighting}[]
\NormalTok{inductive clos\_refl\_trans (R : relation X) : relation X}
\NormalTok{  | refl : ∀ a : X,}
\NormalTok{      clos\_refl\_trans R a a}
\NormalTok{  | step : ∀ a b : X,}
\NormalTok{      R a b → clos\_refl\_trans R a b}
\NormalTok{  | trans : ∀ a b c : X,}
\NormalTok{      clos\_refl\_trans R a b → clos\_refl\_trans R b c →}
\NormalTok{      clos\_refl\_trans R a c}
\end{Highlighting}
\end{Shaded}

\textbf{Lema 1.1.1.} La clausura reflexiva-transitiva de toda relación
es igual a la clausura transitiva de su clausura reflexiva.
\[ R^* = (R^{=})^+ \]

\begin{proof}
\leavevmode

\begin{itemize}
\item
  Para demostrar \(R^* \subseteq(R^{=})^+\), procedemos por inducción
  sobre \(R^*(a,b)\).

  \begin{itemize}
  \item
    En el caso de la regla \textsc{refl}, tenemos \(a=b\): obtenemos
    \(R^{=}(a,a)\) mediante \textsc{refl} y \((R^{=})^+(a,a)\) mediante
    \textsc{step}.
  \item
    En el caso de \textsc{step}, tenemos \(R(a,b)\), que incorporamos
    primero a \(R^{=}\) y después a \((R^{=})^+\) mediante las
    respectivas reglas \textsc{step}.
  \item
    Finalmente, en el caso de \textsc{trans}, las hipótesis de inducción
    proporcionan \((R^{=})^+(a,b)\) y \((R^{=})^+(b,c)\), de donde
    obtenemos \((R^{=})^+(a,c)\) mediante \textsc{trans}.
  \end{itemize}
\item
  Para demostrar \((R^{=})^+\subseteq R^*\), procedemos por inducción
  sobre \((R^{=})^+(a,b)\).

  \begin{itemize}
  \item
    En el caso de \textsc{step}, tenemos \(R^{=}(a,b)\). Según se haya
    utilizado \textsc{refl} o \textsc{step} de \(R^{=}\), hay dos
    posibilidades: \(a=b\) o \(R(a,b)\). Aplicando, respectivamente,
    \textsc{refl} o \textsc{step} de \(R^*\), obtenemos \(R^*(a,b)\).
  \item
    Finalmente, en el caso de \textsc{trans}, las hipótesis de inducción
    proporcionan \(R^*(a,b)\) y \(R^*(b,c)\), de donde obtenemos
    \(R^*(a,c)\) mediante \textsc{trans}.
  \end{itemize}
\end{itemize}

\end{proof}

Su enunciado en Lean es:

\begin{Shaded}
\begin{Highlighting}[]
\KeywordTok{lemma}\NormalTok{ clos\_refl\_trans\_eq : ∀ R : relation X, }
\NormalTok{    clos\_refl\_trans R = clos\_trans (clos\_refl R)}
\end{Highlighting}
\end{Shaded}

\section{1.2. Propiedad del diamante y
confluencia}\label{propiedad-del-diamante-y-confluencia}

Una relación \(R\) tiene la \emph{propiedad del diamante} si y solo si
\[
\forall m,m_1,m_2,\quad
R(m,m_1)\implies R(m,m_2)\implies
\exists m_3,\;
\bigl[R(m_1,m_3)\land R(m_2,m_3)\bigr].
\]

En Lean:

\begin{Shaded}
\begin{Highlighting}[]
\KeywordTok{def}\NormalTok{ diamond (R : relation X) : }\KeywordTok{Prop}\NormalTok{ :=}
\NormalTok{  ∀ m m₁ m₂, R m m₁ → R m m₂ →}
\NormalTok{    ∃ m₃, R m₁ m₃ ∧ R m₂ m₃}
\end{Highlighting}
\end{Shaded}

Escribiendo \(a \longrightarrow b\) en lugar de \(R(a,b)\), gráficamente
esto significa que siempre podemos completar el siguiente diagrama:

\begin{center}
\begin{tikzpicture}[>=stealth, scale=0.7]
  \node (m)  at (0, 1.5) {$m$};
  \node (m1) at (-1.5, 0) {$m_1$};
  \node (m2) at (1.5, 0) {$m_2$};
  \node (m3) at (0, -1.5) {$m_3$};

  \draw[->] (m) -- (m1);
  \draw[->] (m) -- (m2);
  \draw[->, dashed] (m1) -- (m3);
  \draw[->, dashed] (m2) -- (m3);
\end{tikzpicture}
\end{center}

Una relación \(R\) es \emph{confluente} si su clausura
reflexiva-transitiva \(R^*\) tiene la propiedad del diamante.

\begin{Shaded}
\begin{Highlighting}[]
\KeywordTok{def}\NormalTok{ confluent (R : relation X) : }\KeywordTok{Prop}\NormalTok{ :=}
\NormalTok{  diamond (clos\_refl\_trans R)}
\end{Highlighting}
\end{Shaded}

\textbf{Lema 1.2.1.} Si una relación \(R\) tiene la propiedad del
diamante, entonces su clausura transitiva \(R^+\) también la tiene.

\begin{proof}
Sean \(a\), \(b\) y \(c\) tales que \(R^+(a,b)\) y \(R^+(a,c)\). Tenemos
dos sucesiones finitas de uno o más pasos de \(R\) con origen en \(a\).
Colocamos la primera en el borde superior de una cuadrícula y la segunda
en el borde izquierdo.

\begin{figure}[htbp]
\centering
\begin{tikzpicture}[x=1.6cm, y=1.3cm]
  \coordinate (a)  at (0,0);
  \coordinate (a1) at (1,0);
  \coordinate (a2) at (2,0);
  \coordinate (b)  at (3,0);

  \coordinate (c1) at (0,-1);
  \coordinate (u)  at (1,-1);
  \coordinate (v)  at (2,-1);
  \coordinate (w)  at (3,-1);

  \coordinate (c) at (0,-2);
  \coordinate (s) at (1,-2);
  \coordinate (t) at (2,-2);
  \coordinate (d) at (3,-2);

  \draw[->, thick] (a) -- (a1);
  \draw[->, thick] (a1) -- (a2);
  \draw[->, thick] (a2) -- (b);
  \draw[->, thick] (a) -- (c1);
  \draw[->, thick] (c1) -- (c);

  \draw[->, dashed] (c1) -- (u);
  \draw[->, dashed] (u) -- (v);
  \draw[->, dashed] (v) -- (w);

  \draw[->, dashed] (c) -- (s);
  \draw[->, dashed] (s) -- (t);
  \draw[->, dashed] (t) -- (d);

  \draw[->, dashed] (a1) -- (u);
  \draw[->, dashed] (a2) -- (v);
  \draw[->, dashed] (b) -- (w);

  \draw[->, dashed] (u) -- (s);
  \draw[->, dashed] (v) -- (t);
  \draw[->, dashed] (w) -- (d);

  \fill (a)  circle (1.5pt) node[above left] {$a$};
  \fill (a1) circle (1.5pt) node[above] {$a_1$};
  \fill (a2) circle (1.5pt) node[above] {$a_2$};
  \fill (b)  circle (1.5pt) node[above right] {$b$};

  \fill (c1) circle (1.5pt) node[left] {$c_1$};
  \fill (u)  circle (1.5pt);
  \fill (v)  circle (1.5pt);
  \fill (w)  circle (1.5pt);

  \fill (c) circle (1.5pt) node[below left] {$c$};
  \fill (s) circle (1.5pt);
  \fill (t) circle (1.5pt);
  \fill (d) circle (1.5pt) node[below right] {$d$};
\end{tikzpicture}
\end{figure}

En cada casilla, los bordes superior e izquierdo son dos pasos de \(R\)
con el mismo origen. Aplicando la propiedad del diamante de \(R\),
obtenemos los dos pasos restantes.

Repetimos este procedimiento de izquierda a derecha y de arriba abajo,
hasta completar la cuadrícula. De esta manera, los bordes derecho e
inferior proporcionan sucesiones de uno o más pasos desde \(b\) y desde
\(c\) hasta el elemento \(d\) situado en la esquina inferior derecha.
Por tanto, \(R^+(b,d)\) y \(R^+(c,d)\).
\end{proof}

El enunciado correspondiente en Lean es:

\begin{Shaded}
\begin{Highlighting}[]
\KeywordTok{lemma}\NormalTok{ diamond\_clos\_trans \{R : relation X\} :}
\NormalTok{    diamond R →}
\NormalTok{    diamond (clos\_trans R)}
\end{Highlighting}
\end{Shaded}

\bookmarksetup{startatroot}

\chapter{\texorpdfstring{Cálculo
\(\lambda\)}{Cálculo \textbackslash lambda}}\label{cuxe1lculo-lambda}

Recordemos la sintaxis habitual de los términos del cálculo \(\lambda\):
\[
M ::= x \mid M\,M \mid \lambda x.\,M.
\] Esta forma de escribir los términos se conoce como
\emph{representación con nombres}.

El uso de nombres requiere cuidado al reemplazar variables libres por
términos. Por ejemplo, para reemplazar la variable libre \(y\) por \(x\)
en \(\lambda x.\,y\), debemos evitar que la variable introducida quede
ligada. Por eso, primero podemos escribir el mismo término como
\(\lambda z.\,y\) y reemplazar entonces \(y\) por \(x\) para obtener
\(\lambda z.\,x\).

En papel podemos identificar mentalmente \(\lambda x.\,y\) con
\(\lambda z.\,y\) y continuar el argumento. En Lean, debemos justificar
ese cambio y comprobar que la variable introducida permanece libre.

Una forma de abordar esta dificultad es la representación \emph{locally
nameless}, que utilizaremos en este trabajo. En ella, las variables
libres conservan un nombre, mientras que las ligadas se representan
mediante números naturales que indican a qué abstracción hacen
referencia. El número \(0\) corresponde a la abstracción más cercana que
contiene la variable ligada, el \(1\) a la siguiente hacia afuera, y así
sucesivamente. Estos números se denominan \emph{índices de de Bruijn}.

Conservar nombres para las variables libres permite seguir utilizándolos
en los enunciados y las demostraciones, como hacemos habitualmente en
papel. Elegimos esta representación para facilitar que, en la prueba
transicional, el lector pueda relacionar los pasos del argumento
matemático con los de su formalización en Lean.

La sintaxis de los términos en la representación \emph{locally nameless}
está dada por \[
t ::=
\mathsf{bvar}\,i \mid
\mathsf{fvar}\,x \mid
\mathsf{abs}\,t \mid
\mathsf{app}\,t\,t,
\] donde \(i\in\mathbb{N}\) es un índice de de Bruijn y \(x\) es el
nombre de una variable libre. En Lean tomamos los nombres como cadenas
de caracteres:

\begin{Shaded}
\begin{Highlighting}[]
\NormalTok{inductive trm : }\KeywordTok{Type}
\NormalTok{  | bvar : ℕ → trm}
\NormalTok{  | fvar : String → trm}
\NormalTok{  | abs  : trm → trm}
\NormalTok{  | app  : trm → trm → trm}
\end{Highlighting}
\end{Shaded}

Para simplificar la escritura, usamos cadenas de caracteres para
variables libres y naturales para variables ligadas:

\begin{Shaded}
\begin{Highlighting}[]
\NormalTok{instance : Coe String trm where}
\NormalTok{  coe := fvar}

\NormalTok{instance (n : ℕ) : OfNat trm n where}
\NormalTok{  ofNat := bvar n}
\end{Highlighting}
\end{Shaded}

Por ejemplo, el término \(\lambda x. \, x \, y\) puede escribirse en
Lean como \texttt{abs\ (app\ 0\ "y")}.

\subsection{Variables libres}\label{variables-libres}

El conjunto de variables libres de un término \(t\), denotado por
\(\operatorname{fv} (t)\), se define recursivamente mediante \[
\begin{array}{l c l}
\operatorname{fv}(\mathsf{bvar}\, i) & \equiv & \emptyset \\ 
\operatorname{fv}(\mathsf{fvar}\, x) & \equiv & \{x\} \\ 
\operatorname{fv}(\mathsf{app}\, t_1\, t_2) & \equiv & \operatorname{fv}(t_1) \cup \operatorname{fv}(t_2) \\ 
\operatorname{fv}(\mathsf{abs}\, t) & \equiv & \operatorname{fv}(t)
\end{array}
\] Su implementación en Lean es:

\begin{Shaded}
\begin{Highlighting}[]
\KeywordTok{def}\NormalTok{ fv (t : trm) : Finset String :=}
  \KeywordTok{match}\NormalTok{ t }\KeywordTok{with}
\NormalTok{  | bvar \_    =\textgreater{} ∅}
\NormalTok{  | fvar x    =\textgreater{} \{x\}}
\NormalTok{  | app t₁ t₂ =\textgreater{} fv t₁ ∪ fv t₂}
\NormalTok{  | abs u     =\textgreater{} fv u}
\end{Highlighting}
\end{Shaded}

\textbf{Lema 2.0.1.} Para todo conjunto finito de nombres \(L\), existe
\(x\notin L\).

\begin{proof}
Hay infinitas cadenas de caracteres y, por tanto, infinitos nombres.
Como \(L\) es finito, no puede contenerlos todos. Por ello, existe
\(x\notin L\).\footnote{Lean permite trabajar con o sin el axioma de
  elección. Mathlib lo utiliza habitualmente, como ocurre en muchas
  demostraciones de la matemática clásica, y la demostración formal que
  empleamos para este lema depende de él.}
\end{proof}

Su enunciado en Lean es:

\begin{Shaded}
\begin{Highlighting}[]
\KeywordTok{lemma}\NormalTok{ exists\_fresh : ∀ L : Finset String,}
\NormalTok{    ∃ x, x ∉ L}
\end{Highlighting}
\end{Shaded}

En las demostraciones de los lemas y del teorema de confluencia
necesitaremos elegir nombres que no aparezcan libres en ciertos términos
ni pertenezcan a determinados conjuntos finitos. Como los nombres que
debemos evitar forman un conjunto finito, el lema anterior garantiza que
podemos elegir uno distinto de ellos. Nos referiremos a esta elección
diciendo que tomamos \emph{un nombre nuevo}.

\section{2.1. Apertura}\label{apertura}

En la representación \emph{locally nameless}, una abstracción tiene la
forma \(\mathsf{abs} \, t\), donde el cuerpo \(t\) utiliza índices de de
Bruijn para referirse a las variables ligadas. Para trabajar con este
cuerpo resulta útil reemplazar las referencias a la abstracción exterior
por una variable libre. A esta operación la llamaremos \emph{apertura}.

Para definirla, utilizamos una operación auxiliar
\(\{k\rightsquigarrow x\}\,t\), donde \(k\) indica el índice que se
reemplaza por la variable libre de nombre \(x\). Al entrar en una
abstracción, el parámetro \(k\) aumenta en uno. La operación se define
recursivamente por

\[
\begin{array}{l c l}
\{k \rightsquigarrow x\} \, (\mathsf{bvar}\, i) & \equiv & \text{si } (i = k) \text{ entonces } (\mathsf{fvar}\, x) \text{ si no } (\mathsf{bvar}\, i) \\ 
\{k \rightsquigarrow x\} \, (\mathsf{fvar}\, y) & \equiv & \mathsf{fvar}\, y \\ 
\{k \rightsquigarrow x\} \, (\mathsf{app}\, t_1 \, t_2) & \equiv & \mathsf{app}\, (\{k \rightsquigarrow x \} \, t_1) \, (\{k \rightsquigarrow x\} \, t_2) \\ 
\{k \rightsquigarrow x\} \, (\mathsf{abs}\, t) & \equiv & \mathsf{abs}\, (\{(k + 1) \rightsquigarrow x\} \, t)
\end{array}
\]

En Lean:

\begin{Shaded}
\begin{Highlighting}[]
\KeywordTok{def}\NormalTok{ open\_rec (k : ℕ) (x : String) (t : trm) : trm :=}
  \KeywordTok{match}\NormalTok{ t }\KeywordTok{with}
\NormalTok{  | bvar i    =\textgreater{} if i = k then fvar x else bvar i}
\NormalTok{  | fvar y    =\textgreater{} fvar y}
\NormalTok{  | app t₁ t₂ =\textgreater{} app (open\_rec k x t₁) (open\_rec k x t₂)}
\NormalTok{  | abs t     =\textgreater{} abs (open\_rec (k + 1) x t)}

\KeywordTok{notation} \StringTok{"\{"}\NormalTok{ k }\StringTok{" \textasciitilde{}\textgreater{} "}\NormalTok{ x }\StringTok{"\} "}\NormalTok{ t =\textgreater{} open\_rec k x t}
\end{Highlighting}
\end{Shaded}

La apertura del cuerpo \(t\) con una variable libre \(x\) se define
tomando \(k=0\): \[
t^x \equiv \{ 0 \rightsquigarrow x \} \, t
\]

En Lean:

\begin{Shaded}
\begin{Highlighting}[]
\KeywordTok{def}\NormalTok{ open\_var (t : trm) (x : String) : trm :=}
\NormalTok{  \{0 \textasciitilde{}\textgreater{} x\} t}

\NormalTok{infixl:80 }\StringTok{" \^{} "}\NormalTok{ =\textgreater{} open\_var}
\end{Highlighting}
\end{Shaded}

Así, \(t^x\) representa el resultado de aplicar la apertura al cuerpo de
una abstracción \(\mathsf{abs} \, t\), utilizando la variable libre
\(x\).

Por ejemplo, consideremos en Lean la abstracción \texttt{abs\ t}, con
\texttt{t\ =\ app\ (abs\ (app\ 0\ 1))\ 0}. Al abrir su cuerpo con el
nombre \texttt{"x"}, obtenemos:

\begin{Shaded}
\begin{Highlighting}[]
\KeywordTok{example}\NormalTok{ :}
\NormalTok{  (app (abs (app 0 1)) 0) \^{} }\StringTok{"x"}\NormalTok{ = (app (abs (app 0 }\StringTok{"x"}\NormalTok{)) }\StringTok{"x"}\NormalTok{) := }\KeywordTok{rfl}
\end{Highlighting}
\end{Shaded}

El índice \texttt{0} que aparece dentro de la abstracción interior
permanece sin cambios, pues hace referencia a dicha abstracción. En
cambio, el índice \texttt{1} que aparece dentro de ella y el índice
\texttt{0} exterior hacen referencia a la abstracción exterior y son
reemplazados por \texttt{"x"}.

\section{2.2. Clausura local}\label{clausura-local}

El tipo \texttt{trm} admite expresiones cuyos índices de de Bruijn no
hacen referencia a ninguna abstracción, como ocurre en \texttt{bvar\ 0}
y en \texttt{abs\ 1}. En el primer caso, el índice no se encuentra bajo
ninguna abstracción y, en el segundo, no corresponde a la única
abstracción presente, mientras que \texttt{abs\ 0} sí representa el
término identidad \(\lambda x.\,x\).

Diremos que un término es localmente cerrado si todos los índices de de
Bruijn que aparecen en él se encuentran ligados por alguna abstracción.
La proposición \(\mathsf{lc} \, t\) expresa que el término \(t\) es
localmente cerrado. Este predicado se define mediante las reglas

\[
\inferrule*[right=\textsc{lc-var}]
  { }
  {\mathsf{lc} \, (\mathsf{fvar}\,x)}
\qquad
\inferrule*[right=\textsc{lc-app}]
  {\mathsf{lc} \, t_1 \\ \mathsf{lc} \, t_2}
  {\mathsf{lc} \, (\mathsf{app} \, t_1 \, t_2)}
\qquad
\inferrule*[right=\textsc{lc-abs}]
  {\forall x \notin L, \quad \mathsf{lc} \, (t^x)}
  {\mathsf{lc} \, (\mathsf{abs}\,t)}
\]

donde \(L\) es algún conjunto finito de nombres.

En Lean, esta definición se expresa de la siguiente manera:

\begin{Shaded}
\begin{Highlighting}[]
\NormalTok{inductive lc : trm → }\KeywordTok{Prop}
\NormalTok{  | lc\_var (x : String) :}
\NormalTok{      lc (fvar x)}
\NormalTok{  | lc\_app \{t₁ t₂ : trm\} :}
\NormalTok{      lc t₁ → lc t₂ → lc (app t₁ t₂)}
\NormalTok{  | lc\_abs (L : Finset String) (t : trm) :}
\NormalTok{      (∀ x ∉ L, lc (t \^{} x)) →}
\NormalTok{      lc (abs t)}
\end{Highlighting}
\end{Shaded}

El constructor \texttt{lc\_var} establece que toda variable libre es
localmente cerrada. El constructor \texttt{lc\_app} indica que una
aplicación es localmente cerrada cuando lo son sus dos componentes. Por
último, \texttt{lc\_abs} emplea \emph{cuantificación cofinita}: exige
elegir un conjunto finito \texttt{L} y demostrar que \texttt{t\ \^{}\ x}
es localmente cerrado para todo nombre \texttt{x} que no pertenezca a
\texttt{L}. La definición no incluye un constructor para términos de la
forma \texttt{bvar\ i}, pues un índice aislado no está ligado por
ninguna abstracción y, por tanto, no es localmente cerrado.

\textbf{Lema 2.2.1.} La apertura no modifica los términos localmente
cerrados. Es decir, si \(t\) es localmente cerrado, entonces
\(\{k\rightsquigarrow x\}\,t=t\) para todo índice \(k\) y todo nombre
\(x\).

Esto se debe a que todos los índices de \(t\) hacen referencia a
abstracciones contenidas en el propio término, mientras que la apertura
solo reemplaza referencias a abstracciones exteriores, por lo que deja
\(t\) sin cambios.

Su enunciado en Lean es:

\begin{Shaded}
\begin{Highlighting}[]
\KeywordTok{lemma}\NormalTok{ open\_rec\_lc \{t : trm\} :}
\NormalTok{    lc t →}
\NormalTok{    ∀ k x, (\{k \textasciitilde{}\textgreater{} x\} t) = t}
\end{Highlighting}
\end{Shaded}

En particular, si \(\mathsf{lc}\,t\), entonces \(t^x=t\).

\begin{Shaded}
\begin{Highlighting}[]
\KeywordTok{lemma}\NormalTok{ open\_var\_lc \{t : trm\} \{x : String\} :}
\NormalTok{    lc t → t \^{} x = t }
\end{Highlighting}
\end{Shaded}

\section{2.3. Cierre}\label{cierre}

Mientras que la apertura reemplaza determinadas referencias ligadas por
una variable libre, el cierre reemplaza las apariciones libres de un
nombre por índices de de Bruijn. En particular, dado un término \(t\) y
un nombre \(x\), permite construir el cuerpo de una abstracción que liga
esas apariciones de \(x\).

Para definirla, utilizamos una operación auxiliar
\(\{k \leftsquigarrow x \} \, t\), donde \(k\) indica el índice de de
Bruijn que debe asignarse a las variables libres con nombre \(x\). Se
define recursivamente por

\[
\begin{array}{l c l}
\{k \leftsquigarrow x\} \, (\mathsf{bvar}\, i) & \equiv & \mathsf{bvar}\, i \\ 
\{k \leftsquigarrow x\} \, (\mathsf{fvar}\, y) & \equiv & \text{si } (x = y) \text{ entonces } (\mathsf{bvar}\, k) \text{ si no } (\mathsf{fvar}\, y) \\ 
\{k \leftsquigarrow x\} \, (\mathsf{app}\, t_1 \, t_2) & \equiv & \mathsf{app}\, (\{k \leftsquigarrow x \} \, t_1) \, (\{k \leftsquigarrow x\} \, t_2) \\ 
\{k \leftsquigarrow x\} \, (\mathsf{abs}\, t) & \equiv & \mathsf{abs}\, (\{(k + 1) \leftsquigarrow x\} \, t)
\end{array}
\]

En Lean:

\begin{Shaded}
\begin{Highlighting}[]
\KeywordTok{def}\NormalTok{ close\_var\_rec (k : ℕ) (x : String) (t : trm) : trm :=}
  \KeywordTok{match}\NormalTok{ t }\KeywordTok{with}
\NormalTok{  | bvar i    =\textgreater{} bvar i}
\NormalTok{  | fvar y    =\textgreater{} if x = y then bvar k else fvar y}
\NormalTok{  | app t₁ t₂ =\textgreater{} app (close\_var\_rec k x t₁) (close\_var\_rec k x t₂)}
\NormalTok{  | abs t     =\textgreater{} abs (close\_var\_rec (k + 1) x t)}

\KeywordTok{notation} \StringTok{"\{"}\NormalTok{ k }\StringTok{" \textless{}\textasciitilde{} "}\NormalTok{ x }\StringTok{"\} "}\NormalTok{ t =\textgreater{} close\_var\_rec k x t}
\end{Highlighting}
\end{Shaded}

El cierre de un término \(t\) respecto de una variable libre \(x\) se
define tomando \(k=0\): \[
t^{\backslash x} \equiv \{0 \leftsquigarrow x \} \, t
\] En Lean:

\begin{Shaded}
\begin{Highlighting}[]
\KeywordTok{def}\NormalTok{ close\_var (t : trm) (x : String) : trm :=}
\NormalTok{  \{0 \textless{}\textasciitilde{} x\} t}

\NormalTok{infixl:80 }\StringTok{" \^{}\textbackslash{}\textbackslash{} "}\NormalTok{ =\textgreater{} close\_var}
\end{Highlighting}
\end{Shaded}

Así, \(t^{\backslash x}\) reemplaza las variables libres con nombre
\(x\) por índices de de Bruijn que, al formar
\(\mathsf{abs} \, (t^{\backslash x})\), hacen referencia a esta
abstracción.

Por ejemplo:

\begin{Shaded}
\begin{Highlighting}[]
\KeywordTok{example}\NormalTok{ :}
\NormalTok{  (app (abs (app 0 }\StringTok{"x"}\NormalTok{)) }\StringTok{"x"}\NormalTok{) \^{}\textbackslash{} }\StringTok{"x"}\NormalTok{ = (app (abs (app 0 1)) 0) := }\KeywordTok{rfl}
\end{Highlighting}
\end{Shaded}

En este ejemplo, la aparición libre de \texttt{"x"} exterior se
reemplaza por \texttt{0}, mientras que la situada dentro de la
abstracción se reemplaza por \texttt{1}.

\textbf{Lema 2.3.1.} Si \(x\notin\operatorname{fv}(t)\), entonces
\((t^x)^{\backslash x}=t\).

Como \(x\) no aparece libre en \(t\), todas las apariciones de
\(\mathsf{fvar}\,x\) en \(t^x\) fueron introducidas por la apertura. Al
cerrar con el mismo nombre, estas apariciones se sustituyen por los
índices que había antes de abrir el término, por lo que se recupera
\(t\).

Su enunciado en Lean es:

\begin{Shaded}
\begin{Highlighting}[]
\KeywordTok{lemma}\NormalTok{ close\_open\_var \{x : String\} \{t : trm\} :}
\NormalTok{    x ∉ fv t →}
\NormalTok{    (t \^{} x) \^{}\textbackslash{} x = t}
\end{Highlighting}
\end{Shaded}

\textbf{Lema 2.3.2.} Si \(x\) no aparece libre en \(t\) ni en \(u\), la
igualdad de sus aperturas \(t^x=u^x\) implica que \(t=u\).

\begin{proof}
Al cerrar ambos lados de \(t^x=u^x\) con \(x\), obtenemos
\((t^x)^{\backslash x}=(u^x)^{\backslash x}\). Por el Lema 2.3.1, el
lado izquierdo es \(t\) y el derecho es \(u\). Por tanto, \(t=u\).
\end{proof}

Su enunciado en Lean es:

\begin{Shaded}
\begin{Highlighting}[]
\KeywordTok{lemma}\NormalTok{ open\_var\_injective \{t u : trm\} \{x : String\} :}
\NormalTok{    x ∉ fv t →}
\NormalTok{    x ∉ fv u →}
\NormalTok{    t \^{} x = u \^{} x →}
\NormalTok{    t = u }
\end{Highlighting}
\end{Shaded}

\section{2.4. Sustitución}\label{sustituciuxf3n}

La sustitución \(t[x:=s]\) reemplaza las apariciones libres del nombre
\(x\) en \(t\) por el término \(s\). En esta representación, sustituir
dentro de una abstracción solo requiere aplicar la sustitución a su
cuerpo, sin elegir ni cambiar nombres ligados. La operación se define
recursivamente por

\[
\begin{array}{l c l}
(\mathsf{bvar}\,i)[x:=s]
  & \equiv & \mathsf{bvar}\,i \\
(\mathsf{fvar}\,y)[x:=s]
  & \equiv &
  \text{si } (x=y) \text{ entonces } s
  \text{ si no } (\mathsf{fvar}\,y) \\
(\mathsf{app}\,t_1\,t_2)[x:=s]
  & \equiv &
  \mathsf{app}\,(t_1[x:=s])\,(t_2[x:=s]) \\
(\mathsf{abs}\,t)[x:=s]
  & \equiv &
  \mathsf{abs}\,(t[x:=s])
\end{array}
\]

En Lean:

\begin{Shaded}
\begin{Highlighting}[]
\KeywordTok{def}\NormalTok{ subst (t : trm) (x : String) (s : trm) : trm :=}
  \KeywordTok{match}\NormalTok{ t }\KeywordTok{with}
\NormalTok{  | bvar i    =\textgreater{} bvar i}
\NormalTok{  | fvar y    =\textgreater{} if x = y then s else fvar y}
\NormalTok{  | app t₁ t₂ =\textgreater{} app (subst t₁ x s) (subst t₂ x s)}
\NormalTok{  | abs t     =\textgreater{} abs (subst t x s)}

\KeywordTok{notation}\NormalTok{ t }\StringTok{"["}\NormalTok{ x }\StringTok{" := "}\NormalTok{ s }\StringTok{"]"}\NormalTok{ =\textgreater{} subst t x s}
\end{Highlighting}
\end{Shaded}

Nos interesa el caso en que sustituimos una variable libre por otra, es
decir, cuando \(s=\mathsf{fvar}\,y\). Esta sustitución conserva la
clausura local: si \(t\) es localmente cerrado, también lo es
\(t[x:=\mathsf{fvar}\,y]\). En efecto, no introduce nuevos índices, y
los que ya aparecen en \(t\) siguen ligados por las mismas
abstracciones.

\textbf{Lema 2.4.1.} Sea \(t\) un término localmente cerrado. Cerrar
\(t\) con \(x\) y abrir el resultado con \(y\) equivale a sustituir las
apariciones libres de \(x\) por \(\mathsf{fvar}\,y\): \[
(t^{\backslash x})^y=t[x:=\mathsf{fvar}\,y].
\]

El cierre reemplaza las apariciones de \(\mathsf{fvar}\,x\) por índices,
que la apertura con \(y\) reemplaza después por \(\mathsf{fvar}\,y\).
Como \(t\) es localmente cerrado, los índices que ya contenía hacen
referencia a abstracciones de su interior y permanecen sin cambios. Así,
el resultado es \(t[x:=\mathsf{fvar}\,y]\).

Su enunciado en Lean es:

\begin{Shaded}
\begin{Highlighting}[]
\KeywordTok{lemma}\NormalTok{ open\_close\_subst \{t : trm\} \{x y : String\} :}
\NormalTok{    lc t → (t \^{}\textbackslash{} x) \^{} y = t[x := fvar y]}
\end{Highlighting}
\end{Shaded}

\bookmarksetup{startatroot}

\chapter{\texorpdfstring{Reducción
\(\eta\)}{Reducción \textbackslash eta}}\label{reducciuxf3n-eta}

Una reducción es, sencillamente, una relación binaria sobre el conjunto
de todos los términos \(\lambda\). La única reducción que nosotros
trataremos es \(\eta\), cuya definición usual es la siguiente: \[
\inferrule*[right=\textsc{eta-red}]
  {x \notin \operatorname{fv}(t)}
  {\lambda x. \, t \, x
   \longrightarrow_{\eta} t}
\qquad
\inferrule*[right=\textsc{eta-app-1}]
  {t_1 \longrightarrow_{\eta} t_1'}
  {t_1 \, t_2
   \longrightarrow_{\eta}
   t_1' \, t_2}
\] \[
\inferrule*[right=\textsc{eta-app-2}]
  {t_2 \longrightarrow_{\eta} t_2'}
  {t_1 \, t_2
   \longrightarrow_{\eta}
   t_1 \, t_2'}
\qquad
\inferrule*[right=\textsc{eta-abs}]
  {t \longrightarrow_{\eta} t'}
  {\lambda x. t
   \longrightarrow_{\eta}
   \lambda x. t'}
\]

Las reglas \textsc{eta-abs}, \textsc{eta-app-1} y \textsc{eta-app-2},
que permiten reducir dentro del cuerpo de una abstracción o en
cualquiera de los componentes de una aplicación, son llamadas
\emph{reglas de congruencia}.

Llamaremos \emph{redex} \(\eta\) a un término de la forma
\(\lambda x.\,t\,x\) con \(x\notin\operatorname{fv}(t)\). Si \(t\) y
\(t'\) son términos tales que \(t \longrightarrow_{\eta} t'\), entonces
decimos que \(t\) se \emph{reduce} a \(t'\) en un solo paso.

Para las clausuras de la reducción \(\eta\), usaremos las siguientes
notaciones habituales:

\begin{enumerate}
\def\labelenumi{\arabic{enumi}.}
\item
  \(\longrightarrow_{\eta}^{=}\) denota la clausura reflexiva de
  \(\longrightarrow_\eta\). Así, \(t \longrightarrow_{\eta}^{=} t'\)
  significa que \(t=t'\) o que \(t \longrightarrow_\eta t'\).
\item
  \(\longrightarrow_{\eta}^{+}\) denota la clausura transitiva de
  \(\longrightarrow_\eta\). Escribimos
  \(t \longrightarrow_{\eta}^{+} t'\) cuando \(t\) se reduce a \(t'\) en
  uno o más pasos.
\item
  \(\twoheadrightarrow_{\eta}\) denota la clausura reflexiva-transitiva
  de \(\longrightarrow_\eta\). Escribimos
  \(t \twoheadrightarrow_\eta t'\) si \(t=t'\) o si una sucesión finita
  de pasos lleva de \(t\) a \(t'\).
\end{enumerate}

Como \(\eta\) es la única reducción que consideraremos, a partir de
ahora omitiremos el subíndice \(\eta\) en la notación de la reducción y
de sus clausuras.

Para formalizar la reducción \(\eta\) mediante la representación
\emph{locally nameless}, debemos adaptar las reglas anteriores.

En la regla \textsc{eta-red}, el término \(\lambda x.\,t\,x\) se
representa como \(\mathsf{abs}\,(\mathsf{app}\,t\,0)\), donde el índice
\(0\) hace referencia a la abstracción exterior. La condición de que
\(t\) sea localmente cerrado garantiza que todos sus índices hacen
referencia a abstracciones contenidas en el propio \(t\) y, en
particular, que ninguno hace referencia a la abstracción exterior que se
elimina. Así se expresa, en esta regla, la condición de que la variable
ligada no aparezca libre en \(t\).

Las reglas se formulan de modo que la reducción relacione únicamente
términos localmente cerrados. Al igual que la definición de
\(\mathsf{lc}\), la regla \textsc{eta-abs} emplea cuantificación
cofinita. Obtenemos así las siguientes reglas:

\[
\inferrule*[right=\textsc{eta-red}]
  {\mathsf{lc} \, t}
  {\mathsf{abs} \, (\mathsf{app} \, t \, 0)
   \longrightarrow t}
\qquad
\inferrule*[right=\textsc{eta-app-1}]
  {\mathsf{lc} \, t_2 \\ 
    t_1 \longrightarrow t_1'}
  {\mathsf{app}\,t_1\,t_2
   \longrightarrow
   \mathsf{app}\,t_1'\,t_2}
\]

\[
\inferrule*[right=\textsc{eta-app-2}]
  {\mathsf{lc} \, t_1 \\
   t_2 \longrightarrow t_2'}
  {\mathsf{app}\,t_1\,t_2
   \longrightarrow
   \mathsf{app} \, t_1 \, t_2'}
\qquad
\inferrule*[right=\textsc{eta-abs}]
  {\forall x\notin L, \quad
   t^x \longrightarrow t'^x}
  {\mathsf{abs}\,t
   \longrightarrow
   \mathsf{abs}\,t'}
\]

En Lean, esta relación se define inductivamente de la siguiente manera:

\begin{Shaded}
\begin{Highlighting}[]
\NormalTok{inductive eta : relation trm}
\NormalTok{  | eta\_red : ∀ t : trm,}
\NormalTok{      lc t →}
\NormalTok{      eta (abs (app t (bvar 0))) t}
\NormalTok{  | eta\_app1 : ∀ t₁ t₁\textquotesingle{} t₂ : trm,}
\NormalTok{      lc t₂ →}
\NormalTok{      eta t₁ t₁\textquotesingle{} →}
\NormalTok{      eta (app t₁ t₂) (app t₁\textquotesingle{} t₂)}
\NormalTok{  | eta\_app2 : ∀ t₁ t₂ t₂\textquotesingle{} : trm,}
\NormalTok{      lc t₁ →}
\NormalTok{      eta t₂ t₂\textquotesingle{} →}
\NormalTok{      eta (app t₁ t₂) (app t₁ t₂\textquotesingle{})}
\NormalTok{  | eta\_abs : ∀ (L : Finset String) (t t\textquotesingle{} : trm),}
\NormalTok{      (∀ x ∉ L, eta (t \^{} x) (t\textquotesingle{} \^{} x)) →}
\NormalTok{      eta (abs t) (abs t\textquotesingle{})}
\end{Highlighting}
\end{Shaded}

\textbf{Lema 3.0.1 (Lema de regularidad).} Si \(r\longrightarrow s\),
entonces \(\mathsf{lc}\,r\) y \(\mathsf{lc}\,s\).

\begin{proof}
Demostramos primero que \(\mathsf{lc}\,r\), por inducción sobre
\(r\longrightarrow s\).

\begin{itemize}
\item
  En el caso de \textsc{eta-red}, la reducción tiene la forma
  \(\mathsf{abs}\,(\mathsf{app}\,t\,0)\longrightarrow t\), con \(t\)
  localmente cerrado. Todos los índices de \(t\) hacen referencia a
  abstracciones contenidas en el propio término, mientras que el índice
  \(0\) que lo acompaña queda ligado por la abstracción exterior. Por
  tanto, \(\mathsf{abs}\,(\mathsf{app}\,t\,0)\) es localmente cerrado.
\item
  En el caso de \textsc{eta-app-1}, la reducción tiene la forma
  \(\mathsf{app}\,t_1\,t_2
  \longrightarrow\mathsf{app}\,t_1'\,t_2\), con
  \(t_1\longrightarrow t_1'\) y \(t_2\) localmente cerrado. Por
  hipótesis de inducción, \(t_1\) es localmente cerrado. Aplicando
  \textsc{lc-app}, concluimos que \(\mathsf{app}\,t_1\,t_2\) es
  localmente cerrado.
\item
  En el caso de \textsc{eta-app-2}, la reducción tiene la forma
  \(\mathsf{app}\,t_1\,t_2
  \longrightarrow\mathsf{app}\,t_1\,t_2'\), con
  \(t_2\longrightarrow t_2'\) y \(t_1\) localmente cerrado. Por
  hipótesis de inducción, \(t_2\) es localmente cerrado. Aplicando
  \textsc{lc-app}, concluimos que \(\mathsf{app}\,t_1\,t_2\) es
  localmente cerrado.
\item
  En el caso de \textsc{eta-abs}, la reducción tiene la forma
  \(\mathsf{abs}\,t\longrightarrow\mathsf{abs}\,t'\), con
  \(t^x\longrightarrow t'^x\) para todo \(x\notin L\), donde \(L\) es un
  conjunto finito de nombres. Por hipótesis de inducción, \(t^x\) es
  localmente cerrado para todo \(x\notin L\). La regla \textsc{lc-abs}
  permite concluir que \(\mathsf{abs}\,t\) es localmente cerrado.
\end{itemize}

La clausura local de \(s\) se demuestra de manera análoga, por inducción
sobre \(r\longrightarrow s\).
\end{proof}

El enunciado correspondiente en Lean es:

\begin{Shaded}
\begin{Highlighting}[]
\KeywordTok{lemma}\NormalTok{ eta\_regular \{r s : trm\} :}
\NormalTok{    eta r s → lc r ∧ lc s}
\end{Highlighting}
\end{Shaded}

\textbf{Lema 3.0.2.} La reducción \(\eta\) no modifica el conjunto de
nombres libres, es decir, si \(r\longrightarrow s\), entonces
\(\operatorname{fv}(r)=\operatorname{fv}(s)\).

En el paso \(\mathsf{abs}\,(\mathsf{app}\,t\,0)\longrightarrow t\), se
elimina la abstracción exterior junto con la aparición de la variable
que liga, representada por el índice \(0\). Como \(t\) es localmente
cerrado, ninguno de sus índices hace referencia a esa abstracción, así,
el paso conserva los nombres libres. Las reglas de congruencia permiten
realizar la reducción dentro de aplicaciones y abstracciones,
conservando esta misma propiedad.

Su enunciado en Lean es:

\begin{Shaded}
\begin{Highlighting}[]
\KeywordTok{lemma}\NormalTok{ eta\_fv\_eq \{r s : trm\} :}
\NormalTok{    eta r s → fv r = fv s}
\end{Highlighting}
\end{Shaded}

\textbf{Lema 3.0.3.} La sustitución de una variable libre por otra
conserva la reducción \(\eta\), es decir, si \(r\longrightarrow s\),
entonces \[
r[x:=\mathsf{fvar}\,y]
\longrightarrow
s[x:=\mathsf{fvar}\,y].
\]

Para entender esta propiedad, consideremos primero la regla
\textsc{eta-red}, que requiere un término de la forma
\(\mathsf{abs}\,(\mathsf{app}\, r \,0)\), con \(r\) localmente cerrado.
La sustitución de una variable libre por otra conserva ambas
condiciones: mantiene la forma del redex y la clausura local de \(r\).
Por ello, la misma regla permite obtener \[
\mathsf{abs}\,
(\mathsf{app}\,(r[x:=\mathsf{fvar}\,y])\,0)
\longrightarrow
r[x:=\mathsf{fvar}\,y].
\]

El mismo razonamiento se aplica cuando el redex se encuentra dentro de
una aplicación o una abstracción, mediante las reglas de congruencia.

Su enunciado en Lean es:

\begin{Shaded}
\begin{Highlighting}[]
\KeywordTok{lemma}\NormalTok{ eta\_subst\_var \{r s : trm\} \{x y : String\} :}
\NormalTok{    eta r s →}
\NormalTok{    eta (r[x := fvar y]) (s[x := fvar y])}
\end{Highlighting}
\end{Shaded}

En la representación con nombres, la regla \textsc{eta-abs} permite
pasar de \(t \longrightarrow t'\) a
\(\lambda x.\, t \longrightarrow\lambda x.\, t'\). Queremos utilizar
esta misma forma de razonar al trabajar con la representación
\emph{locally nameless}.

Para ello, utilizaremos el cierre antes de añadir la abstracción. Por
ejemplo, en \(\mathsf{abs}\,(\mathsf{fvar}\,x)\), la variable \(x\)
sigue libre. Si primero cerramos \(\mathsf{fvar}\,x\) con \(x\),
obtenemos el índice \(0\). Al añadir la abstracción, este índice queda
ligado y obtenemos \(\mathsf{abs}\,0\), que corresponde a
\(\lambda x.\,x\).

Para un término \(t\), esta construcción se escribe
\(\mathsf{abs}\,(t^{\backslash x})\). El siguiente lema permite
conciliar este paso del argumento informal con su formalización en Lean:
a partir de \(t\longrightarrow t'\), obtenemos una reducción entre las
abstracciones construidas de esta manera.

\textbf{Lema 3.0.4.} Para todo nombre \(x\), se cumple: \[
\inferrule*[right=\textsc{abs-close}]
  {t \longrightarrow t'}
  {\mathsf{abs}\,(t^{\backslash x})
    \longrightarrow
    \mathsf{abs}\,(t'^{\backslash x})}
\]

\begin{proof}
Supongamos que \(t\longrightarrow t'\). Por el lema de regularidad,
ambos términos son localmente cerrados. Queremos demostrar que
\(\mathsf{abs}\,(t^{\backslash x})
\longrightarrow
\mathsf{abs}\,(t'^{\backslash x})\).

Aplicaremos \textsc{eta-abs} con \(L=\varnothing\), pues los lemas que
utilizaremos no requieren excluir ningún nombre. Basta, entonces,
demostrar que
\((t^{\backslash x})^y\longrightarrow(t'^{\backslash x})^y\) para todo
nombre \(y\).

Sea \(y\) un nombre cualquiera. El Lema 3.0.3, aplicado a
\(t\longrightarrow t'\), proporciona \[
t[x:=\mathsf{fvar}\,y]
\longrightarrow
t'[x:=\mathsf{fvar}\,y].
\] Como \(t\) y \(t'\) son localmente cerrados, el Lema 2.4.1 nos da las
igualdades \(t[x:=\mathsf{fvar}\,y]=(t^{\backslash x})^y\) y
\(t'[x:=\mathsf{fvar}\,y]=(t'^{\backslash x})^y\). Obtenemos así
\((t^{\backslash x})^y\longrightarrow(t'^{\backslash x})^y\) para todo
nombre \(y\). Esto satisface la premisa de \textsc{eta-abs} con
\(L=\varnothing\). Aplicando esta regla, concluimos que
\(\mathsf{abs}\,(t^{\backslash x})
\longrightarrow
\mathsf{abs}\,(t'^{\backslash x})\).
\end{proof}

Su enunciado en Lean es:

\begin{Shaded}
\begin{Highlighting}[]
\KeywordTok{lemma}\NormalTok{ abs\_close \{t t\textquotesingle{} : trm\} :}
\NormalTok{    ∀ x, eta t t\textquotesingle{} →}
\NormalTok{    eta (abs (t \^{}\textbackslash{} x)) (abs (t\textquotesingle{} \^{}\textbackslash{} x))}
\end{Highlighting}
\end{Shaded}

\section{3.1. Reglas de congruencia para la clausura
reflexiva}\label{reglas-de-congruencia-para-la-clausura-reflexiva}

Las reglas de congruencia de la reducción \(\eta\) se extienden a su
clausura reflexiva. A continuación presentamos las que utilizaremos en
la demostración de la confluencia.

\textbf{Lema 3.1.1.} \[
\inferrule*[right=\textsc{refl-app-1}]
  {\mathsf{lc} \, t_2 \\ 
    t_1 \longrightarrow^{=} t_1'}
  {\mathsf{app}\,t_1\,t_2
   \longrightarrow^{=}
   \mathsf{app}\,t_1'\,t_2}
\]

\begin{proof}
Si \(t_1=t_1'\), las dos aplicaciones son iguales y aplicamos
\textsc{refl}. Si \(t_1\longrightarrow t_1'\), como \(t_2\) es
localmente cerrado, la regla \textsc{eta-app-1} proporciona
\(\mathsf{app} \, t_1 \, t_2 \longrightarrow
\mathsf{app} \, t_1' \, t_2\). Aplicando \textsc{step}, obtenemos la
conclusión.
\end{proof}

Su enunciado en Lean es:

\begin{Shaded}
\begin{Highlighting}[]
\KeywordTok{lemma}\NormalTok{ refl\_app1 \{t₁ t₁\textquotesingle{} t₂ : trm\} :}
\NormalTok{    lc t₂ →}
\NormalTok{    clos\_refl eta t₁ t₁\textquotesingle{} →}
\NormalTok{    clos\_refl eta (app t₁ t₂) (app t₁\textquotesingle{} t₂)}
\end{Highlighting}
\end{Shaded}

\textbf{Lema 3.1.2.} \[
\inferrule*[right=\textsc{refl-app-2}]
  {\mathsf{lc} \, t_1 \\
   t_2 \longrightarrow^{=} t_2'}
  {\mathsf{app}\,t_1\,t_2
   \longrightarrow^{=}
   \mathsf{app} \, t_1 \, t_2'}
\]

\begin{proof}
Si \(t_2=t_2'\), las dos aplicaciones son iguales y aplicamos
\textsc{refl}. Si \(t_2\longrightarrow t_2'\), como \(t_1\) es
localmente cerrado, la regla \textsc{eta-app-2} proporciona
\(\mathsf{app} \, t_1 \, t_2 \longrightarrow
\mathsf{app} \, t_1 \, t_2'\). Aplicando \textsc{step}, obtenemos la
conclusión.
\end{proof}

Su enunciado en Lean es:

\begin{Shaded}
\begin{Highlighting}[]
\KeywordTok{lemma}\NormalTok{ refl\_app2 \{t₁ t₂ t₂\textquotesingle{} : trm\} :}
\NormalTok{    lc t₁ →}
\NormalTok{    clos\_refl eta t₂ t₂\textquotesingle{} →}
\NormalTok{    clos\_refl eta (app t₁ t₂) (app t₁ t₂\textquotesingle{}) }
\end{Highlighting}
\end{Shaded}

\textbf{Lema 3.1.3.} Para todo nombre \(x\), se cumple: \[
\inferrule*[right=\textsc{refl-abs}]
  {t \longrightarrow^{=} t'}
  {\mathsf{abs}\,(t^{\backslash x})
    \longrightarrow^{=}
    \mathsf{abs}\,(t'^{\backslash x})}
\]

\begin{proof}
Supongamos que \(t\longrightarrow^{=}t'\). Si \(t=t'\), las
abstracciones \(\mathsf{abs} \, (t^{\backslash x})\) y
\(\mathsf{abs} \, (t'^{\backslash x})\) son iguales y aplicamos
\textsc{refl}. Si \(t\longrightarrow t'\), el Lema 3.0.4 proporciona una
reducción entre estas abstracciones. Aplicando \textsc{step}, obtenemos
la conclusión.
\end{proof}

Su enunciado en Lean es:

\begin{Shaded}
\begin{Highlighting}[]
\KeywordTok{lemma}\NormalTok{ refl\_abs \{t t\textquotesingle{} : trm\} :}
\NormalTok{    ∀ x, clos\_refl eta t t\textquotesingle{} →}
\NormalTok{    clos\_refl eta (abs (t \^{}\textbackslash{} x)) (abs (t\textquotesingle{} \^{}\textbackslash{} x))}
\end{Highlighting}
\end{Shaded}

\section{3.2. Análisis de las
reducciones}\label{anuxe1lisis-de-las-reducciones}

La forma del término inicial permite identificar las reglas con las que
puede obtenerse una reducción.

Si \(\mathsf{app}\,t_1\,t_2\longrightarrow s\), hay dos posibilidades:

\begin{enumerate}
\def\labelenumi{\arabic{enumi}.}
\item
  Se aplica la regla \textsc{eta-app-1}: \(t_2\) es localmente cerrado y
  existe un término \(t_1'\) tal que \(t_1\longrightarrow t_1'\) y
  \(s=\mathsf{app}\,t_1'\,t_2\).
\item
  Se aplica la regla \textsc{eta-app-2}: \(t_1\) es localmente cerrado y
  existe un término \(t_2'\) tal que \(t_2\longrightarrow t_2'\) y
  \(s=\mathsf{app}\,t_1\,t_2'\).
\end{enumerate}

Si \(\mathsf{abs}\,t\longrightarrow s\), las posibilidades son las
siguientes:

\begin{enumerate}
\def\labelenumi{\arabic{enumi}.}
\item
  Se aplica la regla \textsc{eta-red}: existe un término \(u\)
  localmente cerrado tal que \(t=\mathsf{app}\,u\,0\) y \(s=u\).
\item
  Se aplica la regla \textsc{eta-abs}: existen un conjunto finito de
  nombres \(L\) y un término \(t'\) tales que
  \(t^x\longrightarrow t'^x\) para todo \(x\notin L\) y
  \(s=\mathsf{abs}\,t'\).
\end{enumerate}

En particular, consideremos una reducción
\(\mathsf{abs}\,(\mathsf{app}\,t\,0)\longrightarrow s\), donde \(t\) es
localmente cerrado. Las dos posibilidades para esta reducción se
presentan de la siguiente manera.

Si se aplica \textsc{eta-red}, obtenemos directamente \(s=t\): \[
\inferrule*[right=\textsc{eta-red}]
  {\mathsf{lc}\,t}
  {\mathsf{abs}\,(\mathsf{app}\,t\,0)\longrightarrow t}.
\]

Si se aplica \textsc{eta-abs}, tenemos \(s=\mathsf{abs}\,u\) para algún
término \(u\), y la reducción se obtiene mediante \[
\inferrule*[right=\textsc{eta-abs}]
  {\forall x\notin L,\quad
   (\mathsf{app}\,t\,0)^x\longrightarrow u^x}
  {\mathsf{abs}\,(\mathsf{app}\,t\,0)
   \longrightarrow\mathsf{abs}\,u},
\] donde \(L\) es un conjunto finito de nombres.

En este último caso, el resultado es una abstracción
\(\mathsf{abs}\,u\). El siguiente lema permite determinar la forma de su
cuerpo a partir de la condición que exige la regla \textsc{eta-abs}.

\textbf{Lema 3.2.1.} Sea \(t\) un término localmente cerrado y sea \(L\)
un conjunto finito de nombres. Si
\((\mathsf{app}\,t\,0)^x\longrightarrow u^x\) para todo \(x\notin L\),
entonces existe un término \(t'\) tal que \(t\longrightarrow t'\) y
\(u=\mathsf{app}\,t'\,0\).

\begin{proof}
Elegimos un nombre nuevo \(x\). Como \(x\notin L\), por hipótesis
tenemos \((\mathsf{app}\,t\,0)^x\longrightarrow u^x\). Como \(t\) es
localmente cerrado, al abrir \(\mathsf{app}\,t\,0\) con \(x\) obtenemos
\(\mathsf{app}\,t\,(\mathsf{fvar}\,x)\). Por tanto, \[
\mathsf{app}\,t\,(\mathsf{fvar}\,x)\longrightarrow u^x.
\] Como \(\mathsf{fvar}\,x\) no se reduce, el paso se obtiene mediante
\textsc{eta-app-1}. Por tanto, existe \(t'\) tal que
\(t\longrightarrow t'\) y \[
u^x=\mathsf{app}\,t'\,(\mathsf{fvar}\,x).
\]

Como \(t\longrightarrow t'\), el Lema 3.0.2 garantiza que el nombre
nuevo \(x\) tampoco aparece libre en \(t'\). Cerrando la igualdad
anterior con \(x\), obtenemos \[
(u^x)^{\backslash x}
=
(\mathsf{app}\,t'\,(\mathsf{fvar}\,x))^{\backslash x}.
\] Por el Lema 2.3.1, \((u^x)^{\backslash x}=u\). Además, como \(x\) no
aparece libre en \(t'\), el cierre solo reemplaza \(\mathsf{fvar}\,x\)
por \(0\) en la aplicación. Por tanto, \(u=\mathsf{app}\,t'\,0\).
\end{proof}

Su enunciado en Lean es:

\begin{Shaded}
\begin{Highlighting}[]
\KeywordTok{lemma}\NormalTok{ eta\_redex\_body \{L : Finset String\} \{t u : trm\} :}
\NormalTok{    lc t →}
\NormalTok{    (∀ x ∉ L, eta ((app t (bvar 0)) \^{} x) (u \^{} x)) →}
\NormalTok{    ∃ t\textquotesingle{}, eta t t\textquotesingle{} ∧ u = app t\textquotesingle{} (bvar 0)}
\end{Highlighting}
\end{Shaded}

\bookmarksetup{startatroot}

\chapter{Una prueba transicional}\label{una-prueba-transicional}

\textbf{Teorema 4.0.1.} Si un término \(m\) se reduce en un paso a
\(m_1\) y a \(m_2\), entonces existe un término \(m_3\) tal que
\(m_1\longrightarrow^{=}m_3\) y \(m_2\longrightarrow^{=}m_3\).

\begin{proof}
Sean \(m\), \(m_1\) y \(m_2\) términos tales que
\(m \longrightarrow m_1\) y \(m \longrightarrow m_2\). Procedemos por
inducción sobre \(m\longrightarrow m_1\). En cada caso demostraremos
que, para cualquier término \(m_2\) tal que \(m\longrightarrow m_2\),
existe un término \(m_3\) con \(m_1\longrightarrow^{=}m_3\) y
\(m_2\longrightarrow^{=}m_3\).

\begin{Shaded}
\begin{Highlighting}[]
\KeywordTok{theorem}\NormalTok{ eta\_diamond : ∀ m m₁ m₂,}
\NormalTok{    eta m m₁ → eta m m₂ →}
\NormalTok{    ∃ m₃, clos\_refl eta m₁ m₃ ∧ clos\_refl eta m₂ m₃ := }\KeywordTok{by}
\NormalTok{  intro m m₁ m₂ h₁ h₂}
\NormalTok{  induction h₁ generalizing m₂}
\end{Highlighting}
\end{Shaded}

\textbf{Caso 1.} La reducción \(m\longrightarrow m_1\) tiene la forma \[
\mathsf{app}\,t\,z \longrightarrow \mathsf{app}\,t_1\,z
\] y se obtiene de \(t\longrightarrow t_1\) mediante la regla
\textsc{eta-app-1}, con \(z\) localmente cerrado. La hipótesis de
inducción afirma que, para cualquier término \(v\) tal que
\(t\longrightarrow v\), existe un término \(q\) tal que
\(t_1\longrightarrow^{=} q\) y \(v\longrightarrow^{=} q\).

Distinguimos dos subcasos según la regla utilizada para obtener
\(\mathsf{app}\,t\,z\longrightarrow m_2\).

\begin{Shaded}
\begin{Highlighting}[]
  \CommentTok{{-}{-} Caso 1}
\NormalTok{  case eta\_app1 t t₁ z lc\_z r\_t\_t₁ ih =\textgreater{}}
\NormalTok{    cases h₂}
\end{Highlighting}
\end{Shaded}

\textbf{Caso 1.1.} Tenemos \(m_2=\mathsf{app}\,t_2\,z\), con
\(t\longrightarrow t_2\). Aplicando la hipótesis de inducción al término
\(t_2\) y al paso \(t\longrightarrow t_2\), obtenemos un término \(t_3\)
tal que \(t_1\longrightarrow^{=}t_3\) y \(t_2\longrightarrow^{=}t_3\).

Como \(z\) es localmente cerrado, el Lema 3.1.1, aplicado a
\(t_1\longrightarrow^{=}t_3\) y \(t_2\longrightarrow^{=}t_3\),
proporciona
\(\mathsf{app}\,t_1\,z\longrightarrow^{=}\mathsf{app}\,t_3\,z\) y
\(\mathsf{app}\,t_2\,z\longrightarrow^{=}\mathsf{app}\,t_3\,z\).

Puesto que \(m_1=\mathsf{app}\,t_1\,z\) y \(m_2=\mathsf{app}\,t_2\,z\),
tomamos \(m_3=\mathsf{app}\,t_3\,z\). Obtenemos así
\(m_1\longrightarrow^{=}m_3\) y \(m_2\longrightarrow^{=}m_3\).

\begin{Shaded}
\begin{Highlighting}[]
    \CommentTok{{-}{-} Caso 1.1}
\NormalTok{    case eta\_app1 t₂ r\_t\_t₂ lc\_z =\textgreater{}}
\NormalTok{      rcases ih t₂ r\_t\_t₂ }\KeywordTok{with}\NormalTok{ ⟨t₃, rl\_t₁\_t₃, rl\_t₂\_t₃⟩}
      \KeywordTok{have}\NormalTok{ rl\_app₁ := refl\_app1 lc\_z rl\_t₁\_t₃}
      \KeywordTok{have}\NormalTok{ rl\_app₂ := refl\_app1 lc\_z rl\_t₂\_t₃}
\NormalTok{      use app t₃ z}
\end{Highlighting}
\end{Shaded}

\textbf{Caso 1.2.} Tenemos \(m_2=\mathsf{app} \, t \, z_1\), con
\(z\longrightarrow z_1\) y \(t\) localmente cerrado. Tomamos
\(m_3=\mathsf{app} \, t_1 \, z_1\) y demostramos las dos relaciones
requeridas.

\begin{Shaded}
\begin{Highlighting}[]
    \CommentTok{{-}{-} Caso 1.2}
\NormalTok{    case eta\_app2 z₁ lc\_t r\_z\_z₁ =\textgreater{}}
\NormalTok{      use app t₁ z₁}
\NormalTok{      constructor}
\end{Highlighting}
\end{Shaded}

Como \(t\longrightarrow t_1\), el Lema 3.0.1 garantiza que \(t_1\) es
localmente cerrado. Por el Lema 3.1.2, para obtener
\(\mathsf{app} \, t_1 \, z\longrightarrow^{=} \mathsf{app} \, t_1 \, z_1\),
basta demostrar \(z\longrightarrow^{=}z_1\). La regla \textsc{step},
aplicada a \(z\longrightarrow z_1\), proporciona esta relación.

\begin{Shaded}
\begin{Highlighting}[]
\NormalTok{      · }\KeywordTok{apply}\NormalTok{ refl\_app2 (eta\_regular r\_t\_t₁).2}
        \KeywordTok{exact}\NormalTok{ clos\_refl.step z z₁ r\_z\_z₁}
\end{Highlighting}
\end{Shaded}

Como \(z\longrightarrow z_1\), el Lema 3.0.1 garantiza que \(z_1\) es
localmente cerrado. Por el Lema 3.1.1, para obtener
\(\mathsf{app} \, t \, z_1\longrightarrow^{=} \mathsf{app} \, t_1 \, z_1\),
basta demostrar \(t\longrightarrow^{=}t_1\). La regla \textsc{step},
aplicada a \(t\longrightarrow t_1\), proporciona esta relación.

\begin{Shaded}
\begin{Highlighting}[]
\NormalTok{      · }\KeywordTok{apply}\NormalTok{ refl\_app1 (eta\_regular r\_z\_z₁).2}
        \KeywordTok{exact}\NormalTok{ clos\_refl.step t t₁ r\_t\_t₁}
\end{Highlighting}
\end{Shaded}

\textbf{Caso 2.} La reducción \(m\longrightarrow m_1\) tiene la forma \[
\mathsf{app}\,z\,t \longrightarrow \mathsf{app}\,z\,t_1
\] y se obtiene de \(t\longrightarrow t_1\) mediante la regla
\textsc{eta-app-2}, con \(z\) localmente cerrado. La hipótesis de
inducción afirma que, para cualquier término \(v\) tal que
\(t\longrightarrow v\), existe un término \(q\) tal que
\(t_1\longrightarrow^{=}q\) y \(v\longrightarrow^{=}q\).

Distinguimos dos subcasos según la regla utilizada para obtener
\(\mathsf{app}\,z\,t\longrightarrow m_2\).

\textbf{Caso 2.1.} Tenemos \(m_2=\mathsf{app} \, z_1 \, t\), con
\(z\longrightarrow z_1\) y \(t\) localmente cerrado. Tomamos
\(m_3=\mathsf{app} \, z_1 \, t_1\) y demostramos las dos relaciones
requeridas.

Como \(t\longrightarrow t_1\), el Lema 3.0.1 garantiza que \(t_1\) es
localmente cerrado. Por el Lema 3.1.1, para obtener
\(\mathsf{app} \, z \, t_1\longrightarrow^{=} \mathsf{app} \, z_1 \, t_1\),
basta demostrar \(z\longrightarrow^{=}z_1\). La regla \textsc{step},
aplicada a \(z\longrightarrow z_1\), proporciona esta relación.

Como \(z\longrightarrow z_1\), el Lema 3.0.1 garantiza que \(z_1\) es
localmente cerrado. Por el Lema 3.1.2, para obtener
\(\mathsf{app} \, z_1 \, t\longrightarrow^{=} \mathsf{app} \, z_1 \, t_1\),
basta demostrar \(t\longrightarrow^{=}t_1\). La regla \textsc{step},
aplicada a \(t\longrightarrow t_1\), proporciona esta relación.

\textbf{Caso 2.2.} Tenemos \(m_2=\mathsf{app}\,z\,t_2\), con
\(t\longrightarrow t_2\). Aplicando la hipótesis de inducción al término
\(t_2\) y al paso \(t\longrightarrow t_2\), obtenemos un término \(t_3\)
tal que \(t_1\longrightarrow^{=}t_3\) y \(t_2\longrightarrow^{=}t_3\).

Como \(z\) es localmente cerrado, el Lema 3.1.2, aplicado a
\(t_1\longrightarrow^{=}t_3\) y \(t_2\longrightarrow^{=}t_3\),
proporciona
\(\mathsf{app}\,z\,t_1\longrightarrow^{=}\mathsf{app}\,z\,t_3\) y
\(\mathsf{app}\,z\,t_2\longrightarrow^{=}\mathsf{app}\,z\,t_3\).

Puesto que \(m_1=\mathsf{app}\,z\,t_1\) y \(m_2=\mathsf{app}\,z\,t_2\),
tomamos \(m_3=\mathsf{app}\,z\,t_3\). Obtenemos así
\(m_1\longrightarrow^{=}m_3\) y \(m_2\longrightarrow^{=}m_3\).

\begin{Shaded}
\begin{Highlighting}[]
  \CommentTok{{-}{-} Caso 2}
\NormalTok{  case eta\_app2 z t t₁ lc\_z r\_t\_t₁ ih =\textgreater{}}
\NormalTok{    cases h₂}
    \CommentTok{{-}{-} Caso 2.1}
\NormalTok{    case eta\_app1 z₁ r\_z\_z₁ lc\_t =\textgreater{}}
\NormalTok{      use app z₁ t₁}
\NormalTok{      constructor}
\NormalTok{      · }\KeywordTok{apply}\NormalTok{ refl\_app1 (eta\_regular r\_t\_t₁).2}
        \KeywordTok{exact}\NormalTok{ clos\_refl.step z z₁ r\_z\_z₁}
\NormalTok{      · }\KeywordTok{apply}\NormalTok{ refl\_app2 (eta\_regular r\_z\_z₁).2}
        \KeywordTok{exact}\NormalTok{ clos\_refl.step t t₁ r\_t\_t₁}
    \CommentTok{{-}{-} Caso 2.2}
\NormalTok{    case eta\_app2 t₂ lc\_z r\_t\_t₂ =\textgreater{}}
\NormalTok{      rcases ih t₂ r\_t\_t₂ }\KeywordTok{with}\NormalTok{ ⟨t₃, rl\_t₁\_t₃, rl\_t₂\_t₃⟩}
      \KeywordTok{have}\NormalTok{ rl\_app₁ := refl\_app2 lc\_z rl\_t₁\_t₃}
      \KeywordTok{have}\NormalTok{ rl\_app₂ := refl\_app2 lc\_z rl\_t₂\_t₃}
\NormalTok{      use app z t₃}
\end{Highlighting}
\end{Shaded}

\textbf{Caso 3.} La reducción \(m\longrightarrow m_1\) tiene la forma \[
\mathsf{abs}\,(\mathsf{app}\,t\,0)\longrightarrow t
\] y se obtiene mediante la regla \textsc{eta-red}, con \(t\) localmente
cerrado. Distinguimos dos subcasos según la regla utilizada para obtener
\(\mathsf{abs}\,(\mathsf{app}\,t\,0)\longrightarrow m_2\).

\textbf{Caso 3.1.} La reducción se obtiene mediante la regla
\textsc{eta-red}, por lo que \(m_2=t\).

Tomamos \(m_3=t\). Puesto que \(m_1=t\) y \(m_2=t\), ambas relaciones
requeridas son \(t\longrightarrow^{=}t\). Aplicando la regla
\textsc{refl} al término \(t\), obtenemos \(m_1\longrightarrow^{=}m_3\)
y \(m_2\longrightarrow^{=}m_3\).

\begin{Shaded}
\begin{Highlighting}[]
  \CommentTok{{-}{-} Caso 3}
\NormalTok{  case eta\_red t lc\_t =\textgreater{}}
\NormalTok{    cases h₂}
    \CommentTok{{-}{-} Caso 3.1}
\NormalTok{    case eta\_red lc\_t =\textgreater{}}
\NormalTok{      use t}
\NormalTok{      constructor}
\NormalTok{      · }\KeywordTok{exact}\NormalTok{ clos\_refl.refl t}
\NormalTok{      · }\KeywordTok{exact}\NormalTok{ clos\_refl.refl t}
\end{Highlighting}
\end{Shaded}

\textbf{Caso 3.2.} La reducción se obtiene mediante la regla
\textsc{eta-abs}. Tenemos \(m_2=\mathsf{abs} \, u\) y un conjunto finito
de nombres \(L\) tal que
\((\mathsf{app} \, t \, 0)^x\longrightarrow u^x\) para todo
\(x\notin L\). Como \(t\) es localmente cerrado, el Lema 3.2.1
proporciona un término \(t'\) tal que \(t\longrightarrow t'\) y
\(u=\mathsf{app} \, t' \,0\).

\begin{Shaded}
\begin{Highlighting}[]
    \CommentTok{{-}{-} Caso 3.2}
\NormalTok{    case eta\_abs L u hbody =\textgreater{}}
\NormalTok{      rcases eta\_redex\_body lc\_t hbody }\KeywordTok{with}\NormalTok{ ⟨t\textquotesingle{}, r\_t\_t\textquotesingle{}, hu⟩}
\end{Highlighting}
\end{Shaded}

Usando esta igualdad escribimos
\(m_2=\mathsf{abs} \, (\mathsf{app} \, t' \,0)\). Tomamos \(m_3=t'\).
Debemos demostrar \(t\longrightarrow^{=}t'\) y
\(\mathsf{abs} \, (\mathsf{app} \, t' \, 0) \longrightarrow^{=}t'\).

\begin{Shaded}
\begin{Highlighting}[]
      \KeywordTok{rewrite}\NormalTok{ [hu]}
\NormalTok{      use t\textquotesingle{}}
\NormalTok{      constructor}
\end{Highlighting}
\end{Shaded}

La primera relación se obtiene aplicando la regla \textsc{step} a
\(t\longrightarrow t'\).

Para la segunda, como \(t\longrightarrow t'\), el Lema 3.0.1 garantiza
que \(t'\) es localmente cerrado. La regla \textsc{eta-red} proporciona
entonces \(\mathsf{abs} \, (\mathsf{app} \, t' \,0)\longrightarrow t'\).
Aplicando \textsc{step}, obtenemos
\(\mathsf{abs} \,(\mathsf{app} \, t' \,0) \longrightarrow^{=}t'\).

\begin{Shaded}
\begin{Highlighting}[]
\NormalTok{      · }\KeywordTok{exact}\NormalTok{ clos\_refl.step t t\textquotesingle{} r\_t\_t\textquotesingle{}}
\NormalTok{      · }\KeywordTok{apply}\NormalTok{ clos\_refl.step}
        \KeywordTok{exact}\NormalTok{ eta.eta\_red t\textquotesingle{} (eta\_regular r\_t\_t\textquotesingle{}).2}
\end{Highlighting}
\end{Shaded}

\textbf{Caso 4.} La reducción \(m\longrightarrow m_1\) tiene la forma \[
\mathsf{abs}\,t\longrightarrow\mathsf{abs}\,t_1
\] y se obtiene mediante la regla \textsc{eta-abs}. Existe, por tanto,
un conjunto finito de nombres \(L\) tal que \(t^x\longrightarrow t_1^x\)
para todo \(x\notin L\). La hipótesis de inducción afirma que, para cada
nombre \(x\notin L\) y cualquier término \(v\) con
\(t^x\longrightarrow v\), existe un término \(q\) tal que
\(t_1^x\longrightarrow^{=}q\) y \(v\longrightarrow^{=}q\).

Distinguimos dos subcasos según la regla utilizada para obtener
\(\mathsf{abs}\,t\longrightarrow m_2\).

\begin{Shaded}
\begin{Highlighting}[]
  \CommentTok{{-}{-} Caso 4}
\NormalTok{  case eta\_abs L t t₁ hbody ih =\textgreater{}}
\NormalTok{    cases h₂}
\end{Highlighting}
\end{Shaded}

\textbf{Caso 4.1.} La reducción \(\mathsf{abs}\,t\longrightarrow m_2\)
se obtiene mediante la regla \textsc{eta-red}, de modo que \(m_2\) es
localmente cerrado y \(t=\mathsf{app}\,m_2\,0\). Por tanto, tenemos
\((\mathsf{app}\,m_2\,0)^x\longrightarrow t_1^x\) para todo
\(x\notin L\).

Aplicando el Lema 3.2.1 a la clausura local de \(m_2\) y a esta
condición, obtenemos un término \(t'\) tal que \(m_2\longrightarrow t'\)
y \(t_1=\mathsf{app}\,t'\,0\). En la relación
\(\mathsf{abs}\,t_1\longrightarrow^{=}m_3\) que debemos demostrar,
sustituimos \(t_1\) por \(\mathsf{app}\,t'\,0\). Tomamos \(m_3=t'\). Las
dos relaciones requeridas son entonces
\(\mathsf{abs}\,(\mathsf{app}\,t'\,0)\longrightarrow^{=}t'\) y
\(m_2\longrightarrow^{=}t'\).

Para la primera relación, por la regla \textsc{step} basta demostrar
\(\mathsf{abs}\,(\mathsf{app}\,t'\,0)\longrightarrow t'\). Este paso se
obtiene mediante \textsc{eta-red}, pues el Lema 3.0.1 aplicado a
\(m_2\longrightarrow t'\) establece que \(t'\) es localmente cerrado.

Para la segunda, aplicamos directamente la regla \textsc{step} a
\(m_2\longrightarrow t'\).

\begin{Shaded}
\begin{Highlighting}[]
    \CommentTok{{-}{-} Caso 4.1}
\NormalTok{    case eta\_red lc\_m₂ =\textgreater{}}
\NormalTok{      rcases eta\_redex\_body lc\_m₂ hbody }\KeywordTok{with}\NormalTok{ ⟨t\textquotesingle{}, r\_m₂\_t\textquotesingle{}, ht₁⟩}
      \KeywordTok{rewrite}\NormalTok{ [ht₁]}
\NormalTok{      use t\textquotesingle{}}
\NormalTok{      constructor}
\NormalTok{      · }\KeywordTok{apply}\NormalTok{ clos\_refl.step}
        \KeywordTok{exact}\NormalTok{ eta.eta\_red t\textquotesingle{} (eta\_regular r\_m₂\_t\textquotesingle{}).2}
\NormalTok{      · }\KeywordTok{exact}\NormalTok{ clos\_refl.step m₂ t\textquotesingle{} r\_m₂\_t\textquotesingle{}}
\end{Highlighting}
\end{Shaded}

\textbf{Caso 4.2.} La reducción \(\mathsf{abs} \, t\longrightarrow m_2\)
se obtiene mediante la regla \textsc{eta-abs}. Tenemos
\(m_2=\mathsf{abs} \, t_2\) y un conjunto finito de nombres \(K\) tal
que \(t^y\longrightarrow t_2^y\) para todo \(y\notin K\).

Elegimos un nombre nuevo \(x\).

\begin{Shaded}
\begin{Highlighting}[]
    \CommentTok{{-}{-} Caso 4.2}
\NormalTok{    case eta\_abs K t₂ hbody₂ =\textgreater{}}
\NormalTok{      rcases exists\_fresh (K ∪ L ∪ fv t₁ ∪ fv t₂) }\KeywordTok{with}\NormalTok{ ⟨x, hx⟩}
      \KeywordTok{have}\NormalTok{ hf : x ∉ K ∧ x ∉ L ∧ x ∉ fv t₁ ∧ x ∉ fv t₂ := }\KeywordTok{by}
\NormalTok{        aesop}
\end{Highlighting}
\end{Shaded}

Como \(x\notin K\), la premisa de \textsc{eta-abs} proporciona
\(t^x\longrightarrow t_2^x\). Además, \(x\notin L\), por lo que
aplicamos la hipótesis de inducción al término \(t_2^x\) y al paso
\(t^x\longrightarrow t_2^x\). Obtenemos un término \(t_3\) tal que
\(t_1^x\longrightarrow^{=}t_3\) y \(t_2^x\longrightarrow^{=}t_3\).

\begin{Shaded}
\begin{Highlighting}[]
      \KeywordTok{have}\NormalTok{ r\_t\_t₂ := hbody₂ x hf.1}
\NormalTok{      rcases ih x hf.2.1 (t₂ \^{} x) r\_t\_t₂ }\KeywordTok{with}\NormalTok{ ⟨t₃, rl\_t₁\_t₃, rl\_t₂\_t₃⟩}
\end{Highlighting}
\end{Shaded}

Aplicando el Lema 3.1.3 a \(t_1^x\longrightarrow^{=}t_3\) y
\(t_2^x\longrightarrow^{=}t_3\), obtenemos
\(\mathsf{abs}\,((t_1^x)^{\backslash x})
\longrightarrow^{=}
\mathsf{abs}\,(t_3^{\backslash x})\) y
\(\mathsf{abs}\,((t_2^x)^{\backslash x})
\longrightarrow^{=}
\mathsf{abs}\,(t_3^{\backslash x})\).

Como \(x\) no aparece libre en \(t_1\) ni en \(t_2\), el Lema 2.3.1
proporciona \((t_1^x)^{\backslash x}=t_1\) y
\((t_2^x)^{\backslash x}=t_2\). Usando estas igualdades, obtenemos
\(\mathsf{abs}\,t_1
\longrightarrow^{=}
\mathsf{abs}\,(t_3^{\backslash x})\) y \(\mathsf{abs}\,t_2
\longrightarrow^{=}
\mathsf{abs}\,(t_3^{\backslash x})\).

Puesto que \(m_1=\mathsf{abs} \, t_1\) y \(m_2=\mathsf{abs} \, t_2\),
tomamos \(m_3=\mathsf{abs} \, (t_3^{\backslash x})\).

\begin{Shaded}
\begin{Highlighting}[]
      \KeywordTok{have}\NormalTok{ rl\_abs₁ := refl\_abs x rl\_t₁\_t₃}
      \KeywordTok{have}\NormalTok{ rl\_abs₂ := refl\_abs x rl\_t₂\_t₃}
      \KeywordTok{rewrite}\NormalTok{ [close\_open\_var hf.2.2.1] }\KeywordTok{at}\NormalTok{ rl\_abs₁}
      \KeywordTok{rewrite}\NormalTok{ [close\_open\_var hf.2.2.2] }\KeywordTok{at}\NormalTok{ rl\_abs₂}
\NormalTok{      use abs (t₃ \^{}\textbackslash{} x)}
\end{Highlighting}
\end{Shaded}

Por tanto, si \(m\longrightarrow m_1\) y \(m\longrightarrow m_2\), en
todos los casos hemos obtenido un término \(m_3\) tal que
\(m_1\longrightarrow^{=}m_3\) y \(m_2\longrightarrow^{=}m_3\), como
queríamos demostrar.
\end{proof}

\textbf{Teorema 4.0.2 (Teorema de Church-Rosser).} La reducción \(\eta\)
es confluente.

\begin{proof}
Debemos demostrar que \[
\forall m,m_1,m_2,\quad
m\twoheadrightarrow m_1\implies
m\twoheadrightarrow m_2\implies
\exists m_3,\;
\bigl[
m_1\twoheadrightarrow m_3\land
m_2\twoheadrightarrow m_3
\bigr],
\] es decir, que \(\twoheadrightarrow\) tiene la propiedad del diamante.
Notemos que \(\twoheadrightarrow=(\longrightarrow^{=})^+\) por el Lema
1.1.1. Así, por el Lema 1.2.1, basta demostrar que
\(\longrightarrow^{=}\) tiene la propiedad del diamante.

Sean entonces \(m\), \(m_1\) y \(m_2\) términos tales que
\(m\longrightarrow^{=}m_1\) y \(m\longrightarrow^{=}m_2\). Procedemos
por casos sobre \(m\longrightarrow^{=}m_1\).

\textbf{Caso 1.} La reducción consta de cero pasos, de modo que
\(m_1=m\).

El objetivo queda entonces en encontrar un término \(m_3\) tal que
\(m\longrightarrow^{=}m_3\) y \(m_2\longrightarrow^{=}m_3\). Tomamos
\(m_3=m_2\). La primera relación es la hipótesis
\(m\longrightarrow^{=}m_2\), mientras que la segunda se obtiene
aplicando la regla \textsc{refl} al término \(m_2\).

\textbf{Caso 2.} La reducción de \(m\) a \(m_1\) consta de un paso, es
decir, \(m\longrightarrow m_1\). Procedemos de manera análoga con
\(m\longrightarrow^{=}m_2\).

\textbf{Caso 2.1.} Esta reducción consta de cero pasos, por lo que
\(m_2=m\).

Buscamos, por tanto, un término \(m_3\) tal que
\(m_1\longrightarrow^{=}m_3\) y \(m\longrightarrow^{=}m_3\). Tomamos
\(m_3=m_1\). La primera relación se obtiene aplicando la regla
\textsc{refl} al término \(m_1\), mientras que para la segunda aplicamos
la regla \textsc{step} a \(m\longrightarrow m_1\).

\textbf{Caso 2.2.} El término \(m\) se reduce también a \(m_2\) en un
paso. Aplicando el Teorema 4.0.1 a \(m\longrightarrow m_1\) y
\(m\longrightarrow m_2\), obtenemos un término \(m_3\) tal que
\(m_1\longrightarrow^{=}m_3\) y \(m_2\longrightarrow^{=}m_3\).

En todos los casos hemos obtenido un término \(m_3\) que satisface las
relaciones requeridas. Por tanto, \(\longrightarrow^{=}\) tiene la
propiedad del diamante y, en consecuencia, la reducción \(\eta\) es
confluente.
\end{proof}

\begin{Shaded}
\begin{Highlighting}[]
\KeywordTok{theorem}\NormalTok{ church\_rosser : confluent eta := }\KeywordTok{by}
\NormalTok{  unfold confluent}
  \KeywordTok{rewrite}\NormalTok{ [clos\_refl\_trans\_eq]}
  \KeywordTok{apply}\NormalTok{ diamond\_clos\_trans}
\NormalTok{  intro m m₁ m₂ h₁ h₂}
\NormalTok{  cases h₁}
  \CommentTok{{-}{-} Caso 1}
\NormalTok{  case refl =\textgreater{}}
\NormalTok{    use m₂}
\NormalTok{    constructor}
\NormalTok{    · }\KeywordTok{exact}\NormalTok{ h₂}
\NormalTok{    · }\KeywordTok{exact}\NormalTok{ clos\_refl.refl m₂}
  \CommentTok{{-}{-} Caso 2}
\NormalTok{  case step hstep₁ =\textgreater{}}
\NormalTok{    cases h₂}
    \CommentTok{{-}{-} Caso 2.1}
\NormalTok{    case refl =\textgreater{}}
\NormalTok{      use m₁}
\NormalTok{      constructor}
\NormalTok{      · }\KeywordTok{exact}\NormalTok{ clos\_refl.refl m₁}
\NormalTok{      · }\KeywordTok{exact}\NormalTok{ clos\_refl.step m m₁ hstep₁}
    \CommentTok{{-}{-} Caso 2.2}
\NormalTok{    case step hstep₂ =\textgreater{}}
      \KeywordTok{exact}\NormalTok{ eta\_diamond m m₁ m₂ hstep₁ hstep₂}
\end{Highlighting}
\end{Shaded}

\bookmarksetup{startatroot}

\chapter*{Conclusión}\label{conclusiuxf3n}
\addcontentsline{toc}{chapter}{Conclusión}

\markboth{Conclusión}{Conclusión}

En este trabajo desarrollamos una prueba transicional de la confluencia
de la reducción \(\eta\) en el cálculo \(\lambda\), utilizando Lean para
su componente formal. Tomamos como referencia la demostración presentada
por Barendregt \autocite{barendregt1984} y seguimos la propuesta de
pruebas transicionales de Miranda-Perea y González Huesca
\autocite{miranda2022}. El trabajo consiste en conciliar este argumento
matemático con su formalización, explicando la correspondencia entre sus
pasos, incluyendo las condiciones y los resultados auxiliares
necesarios.

Para la formalización utilizamos la representación \emph{locally
nameless}, siguiendo el trabajo de Charguéraud
\autocite{chargueraud2012}. Definimos en Lean las operaciones de
apertura, cierre y sustitución y demostramos las propiedades que
utilizamos. Entre las construcciones que retomamos de la literatura se
encuentra la de cerrar un término con un nombre y después añadir una
abstracción. En nuestro desarrollo, el Lema 3.0.4 permite obtener de
esta manera una reducción entre abstracciones: \[
\inferrule*[right=\textsc{abs-close}]
  {t \longrightarrow t'}
  {\mathsf{abs}\,(t^{\backslash x})
    \longrightarrow
    \mathsf{abs}\,(t'^{\backslash x})}
\]

Retomamos esta construcción porque permite razonar sobre las
abstracciones de una manera cercana a la representación con nombres. El
lema justifica ese paso en nuestra formalización y reúne los detalles
necesarios para utilizarlo. Así, al aplicarlo en la prueba transicional,
podemos seguir el argumento matemático y reconocer cómo se expresa en
Lean.

Una continuación sería elaborar una prueba transicional de la
confluencia de la reducción \(\beta\). Para ello, podríamos seguir el
método de reducción paralela de Tait y Martin-Löf
\autocite{barendregt1984} y tomar como referencia formalizaciones
existentes, como la de Nipkow \autocite{nipkow2001}.

También podría desarrollarse una prueba transicional de la conmutación
de \(\beta\) y \(\eta\), en el sentido de que: \[
\text{si }
M\longrightarrow_\beta N_1 \land 
M\longrightarrow_\eta N_2
\text{ entonces existe un término $Q$ tal que }
N_1\longrightarrow_\eta Q \land 
N_2\longrightarrow_\beta Q.
\] Esta propiedad, junto con la confluencia de ambas reducciones,
permite concluir la confluencia de \(\beta\eta\).

En ambos casos, el objetivo sería conciliar los argumentos matemáticos
con sus formalizaciones mediante pruebas transicionales que hagan
explícita la correspondencia entre sus pasos.

\bookmarksetup{startatroot}

\chapter*{Bibliografía}\label{bibliografuxeda}
\addcontentsline{toc}{chapter}{Bibliografía}

\markboth{Bibliografía}{Bibliografía}

\printbibliography[heading=none]





\end{document}
